\chapter{Measure and Integration}\label{measure-theory}
\chapterstatus{complete}

\section{Introduction}\label{measure-theory:section-introduction}

The Lebesgue integral extends the Riemann integral to a class of functions that is closed under the limits needed in analysis: the
monotone and dominated convergence theorems (Theorems~\ref{measure-theory:theorem-monotone-convergence}
and~\ref{measure-theory:theorem-dominated-convergence}) and the completeness of the spaces $L^p$
(Theorem~\ref{measure-theory:theorem-riesz-fischer}) are the basis of the functional analysis used to prove the Hodge theorem. This
chapter constructs Lebesgue measure, proves the convergence theorems, Littlewood's three principles (approximation of sets by boxes, Lusin's
theorem, Egorov's theorem), the Lebesgue differentiation theorem with its consequences (the density theorem and the fundamental theorem of calculus for
Lebesgue integrals, with the Cantor function as a counterexample), Fubini's theorem with the layer cake formula and the Hardy--Littlewood maximal theorem, the
change of variables formula, the Radon--Nikodym theorem, and the characterisation of indefinite integrals as the absolutely continuous functions. It ends with an introduction to
geometric measure theory: Hausdorff measures and Hausdorff dimension (with the Cantor set as example), Rademacher's theorem on the differentiability of Lipschitz maps, and the area formula,
which computes the $k$-dimensional measure of parametrised sets (Sections~\ref{measure-theory:section-hausdorff}--\ref{measure-theory:section-area-formula}).
References: \cite{rudin-rca} and \cite{folland-real}.

\noindent\textbf{Prerequisites.} Chapter~\ref{real-analysis} (Real Analysis).

We use the extended half-line $[0, \infty]$ with $a + \infty = \infty$ and $0 \cdot \infty = 0$.

\section{Measures}\label{measure-theory:section-measures}

\begin{definition}\label{measure-theory:definition-measure}
A \emph{$\sigma$-algebra} on a set $X$ is a set $\mathcal{A}$ of subsets containing $X$ and closed under complements and countable unions. The $\sigma$-algebra \emph{generated} by a family of
subsets is the intersection of all $\sigma$-algebras containing it; for a topological space, the \emph{Borel} $\sigma$-algebra $\mathcal{B}(X)$ is generated by the open sets. A \emph{measure} on
$(X, \mathcal{A})$ is a map $\mu \colon \mathcal{A} \to [0, \infty]$ with $\mu(\emptyset) = 0$ and $\mu(\bigsqcup_k A_k) = \sum_k\mu(A_k)$ for countable disjoint families. It is \emph{finite} if
$\mu(X) < \infty$, \emph{$\sigma$-finite} if $X$ is a countable union of sets of finite measure, and \emph{complete} if subsets of measure zero sets are measurable. A property holds
\emph{almost everywhere} (a.e.) if it holds outside a set of measure zero.
\end{definition}

\begin{proposition}\label{measure-theory:proposition-measure-properties}
Let $\mu$ be a measure. If $A \subseteq B$, then $\mu(A) \leq \mu(B)$; $\mu(\bigcup_kA_k) \leq \sum_k\mu(A_k)$; if $A_1 \subseteq A_2 \subseteq \cdots$, then $\mu(\bigcup_kA_k) = \lim_k\mu(A_k)$; and if
$A_1 \supseteq A_2 \supseteq \cdots$ with $\mu(A_1) < \infty$, then $\mu(\bigcap_kA_k) = \lim_k\mu(A_k)$.
\end{proposition}

\begin{proof}
$\mu(B) = \mu(A) + \mu(B \setminus A)$. Replacing $A_k$ by $A_k \setminus \bigcup_{j<k}A_j$ gives disjoint sets with the same union and smaller measures. For increasing sets,
$\bigcup_kA_k = \bigsqcup_k(A_k \setminus A_{k-1})$, and the partial sums of the measures are $\mu(A_k)$. For decreasing sets apply this to $A_1 \setminus A_k$.
\end{proof}

\begin{definition}\label{measure-theory:definition-measurable-function}
A map $f \colon (X, \mathcal{A}) \to (Y, \mathcal{B})$ is \emph{measurable} if $f^{-1}(B) \in \mathcal{A}$ for $B \in \mathcal{B}$. Functions to $[-\infty, \infty]$, $\RR$ or $\CC$ are measurable for the Borel
$\sigma$-algebras. A \emph{simple function} is a finite combination $\sum_ia_i\mathbf{1}_{A_i}$ of indicator functions of measurable sets.
\end{definition}

\begin{proposition}\label{measure-theory:proposition-measurable-functions}
\begin{enumerate}
\item $f \colon X \to [-\infty, \infty]$ is measurable if and only if $\{f > a\}$ is measurable for all $a \in \RR$.
\item If $f_1, \ldots, f_k \colon X \to \RR$ are measurable and $\Phi \colon \RR^k \to \RR$ is continuous, $\Phi(f_1, \ldots, f_k)$ is measurable; so sums and products of measurable functions are
measurable, and $f \colon X \to \CC$ is measurable if and only if its real and imaginary parts are.
\item Pointwise $\sup$, $\inf$, $\limsup$ and $\liminf$ of sequences of measurable functions are measurable.
\item Every measurable $f \colon X \to [0, \infty]$ is the pointwise limit of an increasing sequence of simple functions $0 \leq s_1 \leq s_2 \leq \cdots$.
\end{enumerate}
\end{proposition}

\begin{proof}
(1) The sets $(a, \infty]$ generate the Borel $\sigma$-algebra of $[-\infty, \infty]$, and the sets $B$ with $f^{-1}(B)$ measurable form a $\sigma$-algebra. (2) Open subsets of $\RR^k$ are countable unions of
products of open intervals, and $(f_1, \ldots, f_k)^{-1}$ of such a product is a finite intersection of measurable sets; so $(f_1, \ldots, f_k)$ is measurable for Borel sets, and so is its composite with $\Phi$.
(3) $\{\sup_nf_n > a\} = \bigcup_n\{f_n > a\}$, $\inf_nf_n = -\sup_n(-f_n)$, and $\limsup f_n = \inf_n\sup_{k \geq n}f_k$. (4) Let $s_n = \min(n, 2^{-n}\lfloor 2^nf \rfloor)$.
\end{proof}

\section{Outer measures and Lebesgue measure}\label{measure-theory:section-lebesgue}

\begin{definition}\label{measure-theory:definition-outer-measure}
An \emph{outer measure} on $X$ is a map $\mu^* \colon \mathcal{P}(X) \to [0, \infty]$ with $\mu^*(\emptyset) = 0$, monotone and countably subadditive. A set $A$ is \emph{$\mu^*$-measurable} if
$\mu^*(E) = \mu^*(E \cap A) + \mu^*(E \setminus A)$ for all $E \subseteq X$.
\end{definition}

\begin{theorem}[Carath\'eodory]\label{measure-theory:theorem-caratheodory}
The $\mu^*$-measurable sets form a $\sigma$-algebra $\mathcal{M}$, and $\mu^*|_\mathcal{M}$ is a complete measure.
\end{theorem}

\begin{proof}
By subadditivity, $A$ is measurable if and only if $\mu^*(E \cap A) + \mu^*(E \setminus A) \leq \mu^*(E)$ for all $E$. The definition is symmetric in $A$ and $X \setminus A$. For measurable $A, B$ and any $E$,
\[
\mu^*(E \cap (A \cup B)) + \mu^*(E \setminus (A \cup B)) \leq \mu^*(E \cap A) + \mu^*((E \setminus A) \cap B) + \mu^*((E \setminus A) \setminus B) = \mu^*(E \cap A) + \mu^*(E \setminus A) = \mu^*(E),
\]
so $\mathcal{M}$ is closed under finite unions. Let $A_1, A_2, \ldots \in \mathcal{M}$ be disjoint, $B_n = A_1 \cup \cdots \cup A_n$ and $B = \bigcup_nA_n$. Measurability of $A_n$ gives
$\mu^*(E \cap B_n) = \mu^*(E \cap A_n) + \mu^*(E \cap B_{n-1})$, so $\mu^*(E \cap B_n) = \sum_{j \leq n}\mu^*(E \cap A_j)$, and
$\mu^*(E) = \mu^*(E \cap B_n) + \mu^*(E \setminus B_n) \geq \sum_{j \leq n}\mu^*(E \cap A_j) + \mu^*(E \setminus B)$. Letting $n \to \infty$ and using subadditivity,
$\mu^*(E) \geq \sum_j\mu^*(E \cap A_j) + \mu^*(E \setminus B) \geq \mu^*(E \cap B) + \mu^*(E \setminus B)$. So $B \in \mathcal{M}$, and $E = B$ gives countable additivity. General countable unions reduce
to disjoint ones. If $\mu^*(A) = 0$, then $\mu^*(E \cap A) + \mu^*(E \setminus A) \leq \mu^*(E)$, so $A$ is measurable.
\end{proof}

\begin{definition}\label{measure-theory:definition-lebesgue}
A \emph{box} in $\RR^n$ is a product $Q = \prod_i[a_i, b_i]$ with volume $v(Q) = \prod_i(b_i - a_i)$. The \emph{Lebesgue outer measure} of $A \subseteq \RR^n$ is
$\lambda^*(A) = \inf\{\sum_kv(Q_k) : A \subseteq \bigcup_kQ_k\}$, the infimum over countable covers by boxes. The $\lambda^*$-measurable sets are \emph{Lebesgue measurable}, and $\lambda = \lambda^n$ is
\emph{Lebesgue measure}.
\end{definition}

\begin{theorem}\label{measure-theory:theorem-lebesgue-measure}
$\lambda^*$ is an outer measure, and $\lambda^*(Q) = v(Q)$ for boxes. Every Borel set is Lebesgue measurable, so $\lambda$ is a complete measure on a $\sigma$-algebra containing $\mathcal{B}(\RR^n)$. It is
translation invariant and \emph{regular}: for measurable $E$, $\lambda(E) = \inf\{\lambda(U) : U \supseteq E \text{ open}\} = \sup\{\lambda(K) : K \subseteq E \text{ compact}\}$.
\end{theorem}

\begin{proof}
Countable subadditivity: given covers of $A_k$ with total volume $\leq \lambda^*(A_k) + \varepsilon2^{-k}$, their union covers $\bigcup_kA_k$.

\emph{Volume of a box.} Clearly $\lambda^*(Q) \leq v(Q)$. Let $Q \subseteq \bigcup_kQ_k$ and $\varepsilon > 0$. Enlarge each $Q_k$ to a box $Q'_k$ whose interior contains $Q_k$, with
$v(Q'_k) \leq v(Q_k) + \varepsilon2^{-k}$. By compactness finitely many interiors cover $Q$, say those of $Q'_1, \ldots, Q'_N$. The endpoints of $Q$ and of the $Q'_k$ in each coordinate define a grid
dividing $Q$ into boxes $C$ with disjoint interiors, and $v(Q) = \sum_Cv(C)$ (volume is multiplicative and length is additive). The centre of each $C$ lies in some $Q'_k$, and then $C \subseteq Q'_k$ because
no grid line crosses the interior of $C$. The cells contained in $Q'_k$ have disjoint interiors, so their volumes add up to at most $v(Q'_k)$. Hence $v(Q) \leq \sum_kv(Q'_k) \leq \sum_kv(Q_k) + \varepsilon$.

\emph{Borel sets.} The half-spaces $H = \{x_i \leq c\}$ generate $\mathcal{B}(\RR^n)$, since open sets are countable unions of boxes with rational vertices. Given $E$ and a cover by boxes $Q_k$ with
$\sum v(Q_k) \leq \lambda^*(E) + \varepsilon$, the boxes $Q_k \cap H$ and $Q_k \cap \{x_i \geq c\}$ cover $E \cap H$ and $E \setminus H$ and have total volume $\sum_kv(Q_k)$. So
$\lambda^*(E \cap H) + \lambda^*(E \setminus H) \leq \lambda^*(E) + \varepsilon$, and $H$ is measurable by Theorem~\ref{measure-theory:theorem-caratheodory}.

\emph{Invariance and regularity.} Translation maps boxes to boxes of the same volume. Enlarging the covering boxes slightly into open boxes gives outer regularity. For inner regularity, first let $E$ be bounded, contained in a
closed box $Q$. Given $\varepsilon$, choose open $U \supseteq Q \setminus E$ with $\lambda(U) \leq \lambda(Q \setminus E) + \varepsilon$; then $K = Q \setminus U \subseteq E$ is compact with
$\lambda(E \setminus K) \leq \lambda(U \setminus (Q \setminus E)) \leq \varepsilon$. For general $E$, apply this to $E \cap [-m, m]^n$ and use continuity from below.
\end{proof}

\begin{counterexample}\label{measure-theory:counterexample-vitali}
There are subsets of $\RR$ that are not Lebesgue measurable. Choose (by the axiom of choice) a set $V \subseteq [0, 1]$ containing exactly one element of each class of $\RR/\QQ$. The translates $V + q$, $q \in \QQ \cap [-1, 1]$,
are disjoint and their union lies between $[0, 1]$ and $[-1, 2]$. If $V$ were measurable, translation invariance and countable additivity would give $1 \leq \sum_q\lambda(V) \leq 3$, which is impossible
whether $\lambda(V) = 0$ or $\lambda(V) > 0$.
\end{counterexample}

\begin{lemma}[Dynkin]\label{measure-theory:lemma-dynkin}
Let $\mathcal{P}$ be a family of subsets of $X$ closed under finite intersections, and $\mathcal{L} \supseteq \mathcal{P}$ a family containing $X$, closed under complements and countable disjoint unions.
Then $\mathcal{L}$ contains the $\sigma$-algebra generated by $\mathcal{P}$. Consequently two measures that agree on $\mathcal{P}$ and are finite on an increasing sequence of sets of $\mathcal{P}$ exhausting $X$ agree on
$\sigma(\mathcal{P})$.
\end{lemma}

\begin{proof}
Let $\mathcal{L}_0$ be the smallest such family containing $\mathcal{P}$. For $A \in \mathcal{P}$, the sets $B$ with $A \cap B \in \mathcal{L}_0$ form such a family (for complements use
$A \setminus B = X \setminus ((X \setminus A) \sqcup (A \cap B))$) containing $\mathcal{P}$, hence $\mathcal{L}_0$. Repeating with $A \in \mathcal{L}_0$ shows that $\mathcal{L}_0$ is closed under finite
intersections, hence under all countable unions (disjointify), so it is a $\sigma$-algebra. For the consequence, if $\mu = \nu$ on $\mathcal{P}$ and $\mu(X_m) = \nu(X_m) < \infty$, the sets $B$ with
$\mu(B \cap X_m) = \nu(B \cap X_m)$ form such a family; let $m \to \infty$.
\end{proof}

\section{The integral}\label{measure-theory:section-integral}

\begin{definition}\label{measure-theory:definition-integral}
For a simple function $s = \sum_ia_i\mathbf{1}_{A_i} \geq 0$ with disjoint $A_i$, $\int s\,d\mu = \sum_ia_i\mu(A_i)$ (independent of the representation, by passing to a common refinement). For measurable
$f \colon X \to [0, \infty]$, $\int f\,d\mu = \sup\{\int s\,d\mu : 0 \leq s \leq f \text{ simple}\}$, and $\int_Ef = \int f\mathbf{1}_E$. A measurable $f \colon X \to \CC$ is \emph{integrable} if $\int|f| < \infty$; then
$\int f = \int(\operatorname{Re}f)^+ - \int(\operatorname{Re}f)^- + i(\int(\operatorname{Im}f)^+ - \int(\operatorname{Im}f)^-)$, where $g^\pm = \max(\pm g, 0)$.
\end{definition}

\begin{theorem}[Monotone convergence]\label{measure-theory:theorem-monotone-convergence}
If $0 \leq f_1 \leq f_2 \leq \cdots$ are measurable and $f = \lim f_n$, then $\int f = \lim\int f_n$.
\end{theorem}

\begin{proof}
The integrals increase and are at most $\int f$. Let $0 \leq s \leq f$ be simple and $0 < c < 1$, and $E_n = \{f_n \geq cs\}$, which increase to $X$. Then $\int f_n \geq c\int_{E_n}s$, and
$A \mapsto \int_As$ is a measure, so $\int_{E_n}s \to \int s$ (Proposition~\ref{measure-theory:proposition-measure-properties}). Hence $\lim\int f_n \geq c\int s$ for all such $c$ and $s$.
\end{proof}

\begin{corollary}\label{measure-theory:corollary-integral-properties}
\begin{enumerate}
\item For measurable $f, g \geq 0$ and $a \geq 0$, $\int(af + g) = a\int f + \int g$, and $\int\sum_nf_n = \sum_n\int f_n$.
\item (Fatou) For measurable $f_n \geq 0$, $\int\liminf f_n \leq \liminf\int f_n$.
\item The integral is linear on integrable functions and $|\int f| \leq \int|f|$.
\item For measurable $f \geq 0$, $\int f = 0$ if and only if $f = 0$ a.e.; and $\mu(\{f \geq t\}) \leq \frac1t\int f$ (Chebyshev).
\end{enumerate}
\end{corollary}

\begin{proof}
(1) For simple functions additivity is clear; approximate by increasing simple functions (Proposition~\ref{measure-theory:proposition-measurable-functions}(4)) and use monotone convergence. (2) Apply
monotone convergence to $g_n = \inf_{k \geq n}f_k$, which satisfy $\int g_n \leq \inf_{k \geq n}\int f_k$. (3) Linearity follows from (1) by decomposing into positive and negative parts. For complex $f$, write
$\int f = |\int f|e^{i\theta}$; then $|\int f| = \int\operatorname{Re}(e^{-i\theta}f) \leq \int|f|$. (4) $t\mathbf{1}_{\{f \geq t\}} \leq f$; and $\{f > 0\} = \bigcup_n\{f \geq 1/n\}$.
\end{proof}

\begin{theorem}[Dominated convergence]\label{measure-theory:theorem-dominated-convergence}
Let $f_n$ be measurable with $f_n \to f$ a.e. and $|f_n| \leq g$ a.e. for an integrable $g$. Then $f$ is integrable and $\int|f_n - f| \to 0$; in particular $\int f_n \to \int f$.
\end{theorem}

\begin{proof}
Modifying on a null set, we may assume convergence and domination everywhere. Then $|f| \leq g$, and Fatou's lemma applied to $2g - |f_n - f| \geq 0$ gives
$\int 2g \leq \int 2g - \limsup\int|f_n - f|$.
\end{proof}

\begin{proposition}\label{measure-theory:proposition-riemann-lebesgue}
For continuous $f \colon [a, b] \to \CC$, the Lebesgue integral $\int_{[a,b]}f\,d\lambda$ equals the Riemann integral.
\end{proposition}

\begin{proof}
For a partition $P$ with mesh tending to $0$, the step functions $\sum_i\min_{[t_{i-1}, t_i]}\operatorname{Re}f \cdot \mathbf{1}_{[t_{i-1}, t_i)}$ converge to $\operatorname{Re}f$ uniformly, and their Lebesgue integrals are
the lower sums (Theorem~\ref{measure-theory:theorem-lebesgue-measure}). Apply dominated convergence, and treat the imaginary part in the same way.
\end{proof}

\begin{proposition}\label{measure-theory:proposition-riemann-integrable}
Every bounded Riemann integrable function $f \colon [a, b] \to \RR$ (Definition~\ref{real-analysis:definition-riemann}) is Lebesgue measurable and Lebesgue integrable, with the same integral.
\end{proposition}

\begin{proof}
Choose partitions $P_1, P_2, \ldots$, each refining the previous one, with $U(f, P_k) - L(f, P_k) < 1/k$ (Proposition~\ref{real-analysis:proposition-riemann}(1)). Let $\phi_k$ and $\psi_k$ be the step functions equal, on the interior of each interval of $P_k$,
to the infimum and supremum of $f$ there, and to $f$ at the finitely many partition points. Then $\phi_k \leq f \leq \psi_k$, $\phi_k$ increases and $\psi_k$ decreases (refinement), and their Lebesgue integrals are $L(f, P_k)$ and $U(f, P_k)$,
since points have measure zero. The limits $\phi = \lim\phi_k$ and $\psi = \lim\psi_k$ are measurable, $\phi \leq f \leq \psi$, and by monotone convergence $\int(\psi - \phi) = \lim(U(f, P_k) - L(f, P_k)) = 0$. So $\psi = \phi$ a.e.
(Corollary~\ref{measure-theory:corollary-integral-properties}(4)), hence $f = \phi$ a.e., and $f$ is measurable because Lebesgue measure is complete. Finally $\int f\,d\lambda = \lim\int\phi_k\,d\lambda = \lim L(f, P_k)$, which is the Riemann integral.
\end{proof}

\begin{counterexample}\label{measure-theory:counterexample-dirichlet}
The converse fails. The Dirichlet function $\mathbf{1}_\QQ$ on $[0, 1]$ is Lebesgue integrable with integral $0$, since $\QQ$ is countable and hence of measure zero. It is not Riemann integrable: every nondegenerate interval contains rational
and irrational points, so $L(\mathbf{1}_\QQ, P) = 0$ and $U(\mathbf{1}_\QQ, P) = 1$ for every partition $P$. It is the increasing pointwise limit of the Riemann integrable functions $\mathbf{1}_{\{q_1, \ldots, q_n\}}$, each of integral $0$, for an
enumeration of $\QQ \cap [0, 1]$: the Riemann integral does not behave well under monotone limits, which is what the Lebesgue integral repairs (Theorem~\ref{measure-theory:theorem-monotone-convergence}).
\end{counterexample}

\section{Lebesgue spaces}\label{measure-theory:section-lp}

\begin{definition}\label{measure-theory:definition-lp}
For $1 \leq p < \infty$, $L^p(\mu)$ is the space of measurable $f \colon X \to \CC$ with $\|f\|_p = (\int|f|^p)^{1/p} < \infty$, modulo equality a.e. $L^\infty(\mu)$ consists of the essentially bounded
functions with $\|f\|_\infty = \inf\{c : |f| \leq c \text{ a.e.}\}$, modulo equality a.e. On $L^2(\mu)$ we have the inner product $\langle f, g \rangle = \int f\bar g$.
\end{definition}

\begin{proposition}\label{measure-theory:proposition-holder}
Let $1 \leq p, q \leq \infty$ with $1/p + 1/q = 1$. Then $\int|fg| \leq \|f\|_p\|g\|_q$ (H\"older) and $\|f + g\|_p \leq \|f\|_p + \|g\|_p$ (Minkowski). So $L^p(\mu)$ is a normed space.
\end{proposition}

\begin{proof}
The cases $p = 1$ or $\infty$ are clear. For $1 < p < \infty$ and $a, b > 0$, concavity of $\log$ ($\log'' < 0$) gives $\log(\frac{a^p}{p} + \frac{b^q}{q}) \geq \log a + \log b$, that is, $ab \leq a^p/p + b^q/q$. If
$\|f\|_p = \|g\|_q = 1$, integrating $|fg| \leq |f|^p/p + |g|^q/q$ gives $\int|fg| \leq 1$; the general case follows by scaling. For Minkowski, $|f + g|^p \leq |f||f + g|^{p-1} + |g||f + g|^{p-1}$ and
H\"older with $(p - 1)q = p$ give $\|f + g\|_p^p \leq (\|f\|_p + \|g\|_p)\|f + g\|_p^{p-1}$; also $|f + g|^p \leq 2^p(|f|^p + |g|^p)$ shows $\|f + g\|_p < \infty$. Functions with $\|f\|_p = 0$ vanish a.e.
by Corollary~\ref{measure-theory:corollary-integral-properties}(4).
\end{proof}

\begin{theorem}[Riesz--Fischer]\label{measure-theory:theorem-riesz-fischer}
$L^p(\mu)$ is a Banach space for $1 \leq p \leq \infty$, and every convergent sequence in $L^p$, $p < \infty$, has a subsequence converging a.e.
\end{theorem}

\begin{proof}
Let $p < \infty$ and $(f_n)$ Cauchy. Choose $n_1 < n_2 < \cdots$ with $\|f_{n_{k+1}} - f_{n_k}\|_p < 2^{-k}$, and let $g = \sum_k|f_{n_{k+1}} - f_{n_k}|$. By monotone convergence and Minkowski,
$\|g\|_p \leq 1$, so $g < \infty$ a.e., and $f = f_{n_1} + \sum_k(f_{n_{k+1}} - f_{n_k})$ converges absolutely a.e. Since $|f - f_{n_k}| \leq \sum_{j \geq k}|f_{n_{j+1}} - f_{n_j}|$, Minkowski gives
$\|f - f_{n_k}\|_p \leq 2^{1-k}$, and a Cauchy sequence with a convergent subsequence converges. For $p = \infty$, outside a null set $|f_n - f_m| \leq \|f_n - f_m\|_\infty$ for all $n, m$, so $(f_n)$ converges
uniformly there.
\end{proof}

\begin{proposition}\label{measure-theory:proposition-density}
For $1 \leq p < \infty$, the continuous functions with compact support are dense in $L^p(\RR^n, \lambda)$.
\end{proposition}

\begin{proof}
By Proposition~\ref{measure-theory:proposition-measurable-functions}(4) and dominated convergence, simple functions $\sum a_i\mathbf{1}_{E_i}$ with $\lambda(E_i) < \infty$ are dense. Given $E$ with $\lambda(E) < \infty$ and
$\varepsilon > 0$, choose compact $K \subseteq E \subseteq U$ open with $\lambda(U \setminus K) < \varepsilon$ (Theorem~\ref{measure-theory:theorem-lebesgue-measure}) and, by
Lemma~\ref{general-topology:lemma-lch-shrinking}, a continuous $\phi \colon \RR^n \to [0, 1]$ with compact support in $U$ and $\phi = 1$ on $K$. Then $\|\mathbf{1}_E - \phi\|_p^p \leq \lambda(U \setminus K) < \varepsilon$.
\end{proof}

\section{Modes of convergence}\label{measure-theory:section-modes}

\begin{definition}\label{measure-theory:definition-modes}
Let $f_n$, $f$ be measurable functions on a measure space $(X, \mu)$ with values in $\CC$. Then $f_n \to f$ \emph{almost everywhere} if $f_n(x) \to f(x)$ for $x$ outside a null set; \emph{in measure} if
$\mu(\{|f_n - f| > \varepsilon\}) \to 0$ for every $\varepsilon > 0$; and \emph{in $L^p$} if $\|f_n - f\|_p \to 0$.
\end{definition}

\begin{theorem}[Egorov]\label{measure-theory:theorem-egorov}
Let $\mu(X) < \infty$ and $f_n \to f$ almost everywhere. For every $\delta > 0$ there is a measurable $E$ with $\mu(X \setminus E) < \delta$ such that $f_n \to f$ uniformly on $E$.
\end{theorem}

\begin{proof}
For $k, n \geq 1$ let $E_{k,n} = \bigcup_{m \geq n}\{|f_m - f| \geq 1/k\}$. For fixed $k$ these sets decrease in $n$, and their intersection consists of points where $|f_m - f| \geq 1/k$ for infinitely many $m$, a null set. As
$\mu(X) < \infty$, $\mu(E_{k,n}) \to 0$ as $n \to \infty$ (Proposition~\ref{measure-theory:proposition-measure-properties}). Choose $n_k$ with $\mu(E_{k,n_k}) < \delta2^{-k}$ and let $E = X \setminus \bigcup_kE_{k,n_k}$. Then
$\mu(X \setminus E) < \delta$, and on $E$, $|f_m - f| < 1/k$ for all $m \geq n_k$, for every $k$.
\end{proof}

\begin{proposition}\label{measure-theory:proposition-modes}
\begin{enumerate}
\item If $f_n \to f$ in $L^p$ for some $1 \leq p < \infty$, then $f_n \to f$ in measure.
\item If $f_n \to f$ in measure, a subsequence converges to $f$ almost everywhere.
\item If $\mu(X) < \infty$ and $f_n \to f$ almost everywhere, then $f_n \to f$ in measure.
\end{enumerate}
\end{proposition}

\begin{proof}
(1) By Chebyshev's inequality (Corollary~\ref{measure-theory:corollary-integral-properties}(4)) applied to $|f_n - f|^p$, $\mu(\{|f_n - f| > \varepsilon\}) \leq \varepsilon^{-p}\|f_n - f\|_p^p$. (2) Choose $n_1 < n_2 < \cdots$ with
$\mu(A_k) < 2^{-k}$, where $A_k = \{|f_{n_k} - f| > 2^{-k}\}$. The set $A$ of points lying in infinitely many $A_k$ is contained in $\bigcup_{k \geq K}A_k$ for every $K$, so $\mu(A) \leq \sum_{k \geq K}2^{-k} = 2^{1-K}$ for all $K$, and
$\mu(A) = 0$. For $x \notin A$, $|f_{n_k}(x) - f(x)| \leq 2^{-k}$ for large $k$. (3) Given $\varepsilon, \delta > 0$, Egorov's theorem gives $E$ with $\mu(X \setminus E) < \delta$ and $|f_n - f| < \varepsilon$ on $E$ for large $n$, so
$\mu(\{|f_n - f| > \varepsilon\}) < \delta$ for large $n$.
\end{proof}

\begin{counterexample}\label{measure-theory:counterexample-modes}
No other implications hold in general. Take $X = [0, 1]$ or $\RR$ with Lebesgue measure.
\begin{enumerate}
\item $f_n = n\mathbf{1}_{(0, 1/n)}$ on $[0, 1]$ converges to $0$ everywhere, but $\int f_n = 1$: almost everywhere convergence and convergence in measure do not imply $L^1$ convergence, even on a finite measure space.
\item $f_n = \mathbf{1}_{[n, n+1]}$ on $\RR$ converges to $0$ everywhere but not in measure, as $\lambda(\{f_n > \frac12\}) = 1$: part (3) of Proposition~\ref{measure-theory:proposition-modes} and Egorov's theorem need $\mu(X) < \infty$.
\item (The typewriter sequence.) For $n = 2^k + j$ with $0 \leq j < 2^k$ let $f_n = \mathbf{1}_{[j2^{-k}, (j+1)2^{-k}]}$ on $[0, 1]$. Then $\|f_n\|_1 = 2^{-k} \to 0$, so $f_n \to 0$ in $L^1$ and in measure; but every $x \in [0, 1]$ lies in one of the
intervals for each $k$ and outside most of them, so $f_n(x) = 1$ and $f_n(x) = 0$ both happen infinitely often, and $(f_n(x))$ converges for no $x$. The subsequence in part (2) of
Proposition~\ref{measure-theory:proposition-modes} is necessary.
\end{enumerate}
\end{counterexample}

\section{Littlewood's principles and Lusin's theorem}\label{measure-theory:section-lusin}

Littlewood summarised the theory of Lebesgue measure on $\RR^n$ in three principles: every measurable set of finite measure is nearly a finite union of boxes;
every measurable function is nearly continuous; every convergent sequence of measurable functions is nearly uniformly convergent. The third is Egorov's
theorem (Theorem~\ref{measure-theory:theorem-egorov}); this section makes the first two precise. Throughout, $\lambda$ is Lebesgue measure on $\RR^n$ and
$E \mathbin{\triangle} A = (E \setminus A) \cup (A \setminus E)$ is the symmetric difference.

\begin{lemma}\label{measure-theory:lemma-dyadic-cubes}
Every open set $W \subseteq \RR^n$ is a countable union of closed \emph{dyadic cubes} $\prod_i[k_i2^{-j}, (k_i + 1)2^{-j}]$ ($j \in \NN$, $k_i \in \ZZ$) with pairwise
disjoint interiors, and $\lambda(W)$ is the sum of their volumes.
\end{lemma}

\begin{proof}
Two dyadic cubes of side at most $1$ are either nested or have disjoint interiors, and a dyadic cube of side at most $1$ is contained in only finitely many
dyadic cubes of side at most $1$. Every $x \in W$ lies in a dyadic cube of side at most $1$ contained in $W$, since $W$ contains a ball around $x$. Let
$(Q_k)$ be the dyadic cubes of side at most $1$ contained in $W$ that are maximal for inclusion among these; every dyadic cube of side at most $1$ in $W$ lies in
one of them, so they cover $W$, and distinct maximal cubes have disjoint interiors. There are countably many dyadic cubes. The boundary of a cube is a
finite union of boxes of volume $0$, a null set (Theorem~\ref{measure-theory:theorem-lebesgue-measure}), so the interiors, which are disjoint, have the same
measures as the cubes, and $\lambda(W) = \sum_k\lambda(Q_k)$ by countable additivity.
\end{proof}

\begin{proposition}[First principle]\label{measure-theory:proposition-approximation-boxes}
Let $E \subseteq \RR^n$ be Lebesgue measurable with $\lambda(E) < \infty$, and $\varepsilon > 0$. There is a finite union $A$ of boxes with
$\lambda(E \mathbin{\triangle} A) < \varepsilon$.
\end{proposition}

\begin{proof}
By Theorem~\ref{measure-theory:theorem-lebesgue-measure} there is an open $W \supseteq E$ with $\lambda(W \setminus E) < \varepsilon/2$. Write $W = \bigcup_kQ_k$ as in
Lemma~\ref{measure-theory:lemma-dyadic-cubes}. Since $\sum_k\lambda(Q_k) = \lambda(W) < \infty$, there is $N$ with $\sum_{k > N}\lambda(Q_k) < \varepsilon/2$. Let
$A = Q_1 \cup \cdots \cup Q_N$. Then $\lambda(A \setminus E) \leq \lambda(W \setminus E) < \varepsilon/2$ and
$\lambda(E \setminus A) \leq \lambda(W \setminus A) \leq \sum_{k > N}\lambda(Q_k) < \varepsilon/2$.
\end{proof}

\begin{theorem}[Lusin]\label{measure-theory:theorem-lusin}
Let $E \subseteq \RR^n$ be Lebesgue measurable with $\lambda(E) < \infty$, let $f \colon E \to \CC$ be measurable, and $\varepsilon > 0$. There is a compact set
$K \subseteq E$ with $\lambda(E \setminus K) < \varepsilon$ such that the restriction $f|_K$ is continuous. Moreover there is $g \in C_c(\RR^n)$ with $g = f$ on $K$,
and if $|f| \leq c$ on $E$ we can take $|g| \leq c$.
\end{theorem}

\begin{proof}
\emph{Step 1: indicator functions.} Let $f = \mathbf{1}_A$ with $A \subseteq E$ measurable. By inner regularity
(Theorem~\ref{measure-theory:theorem-lebesgue-measure}) there are compact sets $K_1 \subseteq A$ and $K_2 \subseteq E \setminus A$ with
$\lambda(A \setminus K_1) < \varepsilon/2$ and $\lambda((E \setminus A) \setminus K_2) < \varepsilon/2$. Then $K = K_1 \cup K_2$ is compact with $\lambda(E \setminus K) < \varepsilon$.
As $K_1$ and $K_2$ are disjoint and closed, each is open in $K$, and $f|_K$ is constant on each, hence continuous.

\emph{Step 2: simple functions.} If $f = \sum_{i=1}^ma_i\mathbf{1}_{A_i}$, apply Step 1 to each $A_i$ with $\varepsilon/m$ and intersect the compact sets.

\emph{Step 3: real functions.} By Proposition~\ref{measure-theory:proposition-measurable-functions}(4), applied to $f^+$ and $f^-$, there are simple functions
$s_k$ with $s_k \to f$ pointwise on $E$. By Step 2 there are compact $K_k \subseteq E$ with $\lambda(E \setminus K_k) < \varepsilon2^{-k-2}$ and $s_k|_{K_k}$
continuous. By Egorov's theorem (Theorem~\ref{measure-theory:theorem-egorov}, on the finite measure space $E$) there is a measurable $F \subseteq E$ with
$\lambda(E \setminus F) < \varepsilon/4$ on which $s_k \to f$ uniformly, and by inner regularity a compact $K_0 \subseteq F$ with $\lambda(F \setminus K_0) < \varepsilon/4$.
The compact set $K = K_0 \cap \bigcap_kK_k$ satisfies $\lambda(E \setminus K) < \varepsilon/4 + \varepsilon/4 + \sum_k\varepsilon2^{-k-2} \leq \varepsilon$, and on it the
continuous functions $s_k|_K$ converge uniformly to $f|_K$, which is therefore continuous (Proposition~\ref{general-topology:proposition-uniform-limit}). For
complex $f$, treat the real and imaginary parts with $\varepsilon/2$ each and intersect.

\emph{Extension.} The space $\RR^n$ is metric, hence normal (Proposition~\ref{general-topology:proposition-normal-examples}), and $K$ is closed, so by the Tietze
extension theorem (Theorem~\ref{general-topology:theorem-tietze}) the real and imaginary parts of $f|_K$ extend to continuous functions on $\RR^n$; let $h$ be the
resulting extension of $f|_K$. If $|f| \leq c$, replace $h$ by $r \circ h$, where $r(z) = z$ for $|z| \leq c$ and $r(z) = cz/|z|$ otherwise is a continuous retraction
of $\CC$ onto the closed disc of radius $c$. Finally choose $\phi \in C_c(\RR^n)$ with values in $[0, 1]$ and $\phi = 1$ on $K$
(Lemma~\ref{general-topology:lemma-lch-shrinking}), and let $g = \phi h$.
\end{proof}

\begin{example}\label{measure-theory:example-lusin-dirichlet}
Lusin's theorem says that the \emph{restriction} of $f$ to $K$ is continuous, not that $f$ is continuous at the points of $K$. The Dirichlet function
$\mathbf{1}_\QQ$ on $[0, 1]$ is discontinuous at every point. Enumerate $\QQ \cap [0, 1] = \{q_1, q_2, \ldots\}$ and let $U$ be the union of the intervals
$(q_k - \varepsilon2^{-k-2}, q_k + \varepsilon2^{-k-2})$. Then $\lambda(U) \leq \varepsilon/2$, the set $K = [0, 1] \setminus U$ is compact with
$\lambda([0, 1] \setminus K) < \varepsilon$, and $\mathbf{1}_\QQ$ vanishes on $K$, so its restriction to $K$ is continuous.
\end{example}

\begin{corollary}\label{measure-theory:corollary-lusin-limit}
Every measurable function $f \colon \RR^n \to \CC$ is the almost everywhere limit of a sequence of continuous functions.
\end{corollary}

\begin{proof}
For $k \geq 1$ apply Theorem~\ref{measure-theory:theorem-lusin} to the restriction of $f$ to the ball $B(0, k)$ with $\varepsilon = 2^{-k}$, to get
$g_k \in C_c(\RR^n)$ equal to $f$ on $B(0, k)$ outside a set $N_k$ of measure less than $2^{-k}$. By monotone convergence,
$\int\sum_k\mathbf{1}_{N_k} = \sum_k\lambda(N_k) < \infty$, so $\sum_k\mathbf{1}_{N_k} < \infty$ almost everywhere: almost every $x$ lies in only finitely many $N_k$.
For such $x$ and all large $k > \|x\|$, $g_k(x) = f(x)$.
\end{proof}

\section{Differentiation of integrals}\label{measure-theory:section-differentiation}

For a continuous function $f$ on $\RR$, the fundamental theorem of calculus says that $\frac{1}{2r}\int_{x-r}^{x+r}f \to f(x)$. We show that the averages of
an integrable function over small balls converge to the function almost everywhere. Let $B(x, r)$ be the open Euclidean ball. The map $y \mapsto x + ry$ multiplies
the volume of boxes by $r^n$, hence Lebesgue measure by $r^n$, so $\lambda(B(x, r)) = r^n\lambda(B(0, 1))$. A measurable function is \emph{locally integrable},
$f \in L^1_{\mathrm{loc}}(\RR^n)$, if $\int_K|f| < \infty$ for every compact $K$.

\begin{lemma}[Vitali covering lemma]\label{measure-theory:lemma-vitali-covering}
Every finite family of open balls $B_1, \ldots, B_m$ in $\RR^n$ has a pairwise disjoint subfamily $B_{i_1}, \ldots, B_{i_k}$ with
$\bigcup_iB_i \subseteq \bigcup_j3B_{i_j}$, where $3B(x, r) = B(x, 3r)$. Consequently $\lambda(\bigcup_iB_i) \leq 3^n\sum_j\lambda(B_{i_j})$.
\end{lemma}

\begin{proof}
Choose a ball of largest radius, discard all the balls that meet it, choose a ball of largest radius among the remaining ones, and so on. The chosen balls are
pairwise disjoint. A discarded ball $B = B(c, r)$ meets a chosen ball $B' = B(c', r')$ with $r' \geq r$, namely the one chosen at the step where $B$ was discarded.
If $z \in B \cap B'$ and $y \in B$, then $\|y - c'\| \leq \|y - c\| + \|c - z\| + \|z - c'\| < r + r + r' \leq 3r'$, so $B \subseteq 3B'$. The inequality follows from
$\lambda(3B') = 3^n\lambda(B')$ and subadditivity.
\end{proof}

\begin{definition}\label{measure-theory:definition-maximal-function}
For $f \in L^1_{\mathrm{loc}}(\RR^n)$, $x \in \RR^n$ and $r > 0$, the \emph{average} of $f$ over $B(x, r)$ is
$A_rf(x) = \frac{1}{\lambda(B(x, r))}\int_{B(x, r)}f\,d\lambda$, and the \emph{Hardy--Littlewood maximal function} is $Mf(x) = \sup_{r > 0}A_r|f|(x) \in [0, \infty]$.
\end{definition}

\begin{theorem}[Hardy--Littlewood maximal inequality]\label{measure-theory:theorem-maximal-inequality}
Let $f \in L^1(\RR^n)$ and $t > 0$. The set $\{Mf > t\}$ is open, and $\lambda(\{Mf > t\}) \leq \frac{3^n}{t}\|f\|_1$.
\end{theorem}

\begin{proof}
If $A_r|f|(x) > t$, choose $s > r$ with $(r/s)^nA_r|f|(x) > t$. For $\|x' - x\| < s - r$ we have $B(x', s) \supseteq B(x, r)$, so
$A_s|f|(x') \geq (r/s)^nA_r|f|(x) > t$; hence $\{Mf > t\}$ is open. Let $K \subseteq \{Mf > t\}$ be compact. Every $x \in K$ is the centre of a ball $B_x$ with
$\int_{B_x}|f| > t\lambda(B_x)$, and finitely many of these balls cover $K$. By Lemma~\ref{measure-theory:lemma-vitali-covering} they contain disjoint balls
$B_1, \ldots, B_k$ with
\[
\lambda(K) \leq 3^n\sum_j\lambda(B_j) < \frac{3^n}{t}\sum_j\int_{B_j}|f| \leq \frac{3^n}{t}\|f\|_1.
\]
By inner regularity (Theorem~\ref{measure-theory:theorem-lebesgue-measure}), the same bound holds for $\{Mf > t\}$.
\end{proof}

\begin{theorem}[Lebesgue differentiation theorem]\label{measure-theory:theorem-lebesgue-differentiation}
Let $f \in L^1_{\mathrm{loc}}(\RR^n)$. For almost every $x \in \RR^n$,
\[
\lim_{r \to 0}\frac{1}{\lambda(B(x, r))}\int_{B(x, r)}|f(y) - f(x)|\,d\lambda(y) = 0,
\]
and in particular $A_rf(x) \to f(x)$ as $r \to 0$.
\end{theorem}

\begin{proof}
For $x \in B(0, R)$ and $r < 1$ the quantity only involves $f$ on $B(0, R + 1)$, so, replacing $f$ by $f\mathbf{1}_{B(0, R+1)}$ and letting $R$ run through $\NN$, we
may assume $f \in L^1(\RR^n)$. Let $\Omega_f(x) = \limsup_{r \to 0}A_r|f - f(x)|(x)$. Fix $t > 0$ and $\varepsilon > 0$, and choose $g \in C_c(\RR^n)$ with
$\|f - g\|_1 < \varepsilon$ (Proposition~\ref{measure-theory:proposition-density}). For the continuous $g$, $A_r|g - g(x)|(x) \leq \sup_{B(x, r)}|g - g(x)| \to 0$, so
$\Omega_g = 0$. With $h = f - g$ we have $|f(y) - f(x)| \leq |h(y)| + |h(x)| + |g(y) - g(x)|$, hence $\Omega_f(x) \leq Mh(x) + |h(x)|$, and
\[
\{\Omega_f > 2t\} \subseteq \{Mh > t\} \cup \{|h| > t\}.
\]
By Theorem~\ref{measure-theory:theorem-maximal-inequality} and Chebyshev's inequality (Corollary~\ref{measure-theory:corollary-integral-properties}(4)), the right
side has measure at most $3^n\varepsilon/t + \varepsilon/t$. As $\varepsilon$ is arbitrary, $\{\Omega_f > 2t\}$ has outer measure $0$, so it is a null set (Lebesgue
measure is complete), and so is the union of these sets for $t = 1/k$, $k \geq 1$. Finally $|A_rf(x) - f(x)| \leq A_r|f - f(x)|(x)$.
\end{proof}

The points where the conclusion holds are the \emph{Lebesgue points} of $f$; changing $f$ on a null set does not change the averages, so almost every point is
a Lebesgue point of every representative.

\begin{corollary}[Lebesgue density theorem]\label{measure-theory:corollary-density}
Let $A \subseteq \RR^n$ be Lebesgue measurable. Then $\lambda(A \cap B(x, r))/\lambda(B(x, r)) \to 1$ as $r \to 0$ for almost every $x \in A$, and $\to 0$ for almost
every $x \notin A$.
\end{corollary}

\begin{proof}
Apply Theorem~\ref{measure-theory:theorem-lebesgue-differentiation} to the bounded, hence locally integrable, function $\mathbf{1}_A$, whose average over $B(x, r)$ is
the ratio in question.
\end{proof}

\begin{corollary}[Fundamental theorem of calculus for Lebesgue integrals]\label{measure-theory:corollary-ftc-lebesgue}
Let $f \in L^1([a, b])$ and $F(x) = \int_{[a, x]}f\,d\lambda$. Then $F$ is continuous, and at almost every $x \in (a, b)$ it is differentiable with $F'(x) = f(x)$.
\end{corollary}

\begin{proof}
Extend $f$ by $0$ to $\RR$. For $y \to x$, $|F(y) - F(x)| \leq \int\mathbf{1}_{[x, y]}|f| \to 0$ by dominated convergence, where $[x, y]$ means the interval between
$x$ and $y$. For $h \neq 0$ let $I_h$ be the closed interval between $x$ and $x + h$, of length $|h|$. Then $F(x + h) - F(x) = \pm\int_{I_h}f$, with the sign of
$h$, so in both cases $\frac{F(x + h) - F(x)}{h} - f(x) = \frac{1}{|h|}\int_{I_h}(f - f(x))$. As $I_h \subseteq \overline{B}(x, |h|)$, an interval of length $2|h|$
whose endpoints are null,
\[
\Bigl|\frac{F(x + h) - F(x)}{h} - f(x)\Bigr| \leq \frac{1}{|h|}\int_{B(x, |h|)}|f - f(x)| = 2A_{|h|}|f - f(x)|(x),
\]
which tends to $0$ at every Lebesgue point of $f$.
\end{proof}

\begin{counterexample}\label{measure-theory:counterexample-cantor-function}
\emph{The Cantor function.} A continuous nondecreasing function can have derivative $0$ almost everywhere without being constant, so the formula
$F(b) - F(a) = \int_a^bF'$ fails for such functions. Let $C \subseteq [0, 1]$ be the Cantor set, the set of numbers $\sum_{k \geq 1}a_k3^{-k}$ with all
$a_k \in \{0, 2\}$; these expansions are unique, and $C$ is compact of measure zero, being covered, for each $k$, by $2^k$ intervals of length $3^{-k}$
(Exercise~\ref{measure-theory:exercise-cantor}). Define $c \colon C \to [0, 1]$ by $c(\sum_ka_k3^{-k}) = \sum_k(a_k/2)2^{-k}$.
\begin{enumerate}
\item $c$ is nondecreasing and surjective. If $x < y$ in $C$, their digits first differ at some $k$ with $a_k = 0$ and $b_k = 2$ (a difference $2 \cdot 3^{-k}$ in
position $k$ cannot be compensated by the later digits, which contribute at most $\sum_{j > k}2 \cdot 3^{-j} = 3^{-k}$); then $c(y) - c(x) \geq 2^{-k} - \sum_{j > k}2^{-j} = 0$.
Every $t \in [0, 1]$ has a binary expansion $\sum_kb_k2^{-k}$, and $t = c(\sum_k2b_k3^{-k})$.
\item Let $F(x) = \sup\{c(y) : y \in C, y \leq x\}$ for $x \in [0, 1]$. Then $F$ is nondecreasing, $F = c$ on $C$, and $F([0, 1]) = [0, 1]$. A nondecreasing function
whose image is an interval has no jumps, so $F$ is continuous.
\item The complement $[0, 1] \setminus C$ is a countable union of open intervals $(\alpha, \beta)$ with $\alpha \in C$, removed in the construction; on each of them
$F = c(\alpha)$ is constant, as $C$ has no point in $(\alpha, \beta)$. So $F' = 0$ on $[0, 1] \setminus C$, that is, almost everywhere, while $F(1) - F(0) = 1$.
\end{enumerate}
So $F$ is not the integral of its derivative. The functions that are integrals of their derivatives are exactly the \emph{absolutely continuous} ones
\cite[Chapter 7]{rudin-rca}. Note also that the continuous map $F$ sends the null set $C$ onto $[0, 1]$.
\end{counterexample}

\begin{counterexample}\label{measure-theory:counterexample-non-borel}
\emph{A Lebesgue measurable set that is not a Borel set.} First, every $A \subseteq \RR$ with $\lambda^*(A) > 0$ contains a non-measurable subset. Let $V$ be as in
Counterexample~\ref{measure-theory:counterexample-vitali}. If $S$ is a measurable subset of a translate $V + q$, $q \in \QQ$, then the sets $S - q + q'$,
$q' \in \QQ \cap [-1, 1]$, are pairwise disjoint (they lie in the disjoint translates $V + q'$), contained in $[-1, 2]$, and all of measure $\lambda(S)$; so
$\lambda(S) = 0$. The sets $A \cap (V + q)$, $q \in \QQ$, cover $A$, since every real number differs from an element of $V$ by a rational; if they were all
measurable, $A$ would be a countable union of null sets.

Now let $F$ be the Cantor function and $h(x) = x + F(x)$ on $[0, 1]$. It is continuous and strictly increasing from $[0, 1]$ onto $[0, 2]$, hence a
homeomorphism. It maps each interval $(\alpha, \beta)$ of $[0, 1] \setminus C$ onto an interval of the same length, since $F$ is constant there; these images are
disjoint, so $\lambda(h([0, 1] \setminus C)) = 1$ and $\lambda(h(C)) = 1$. Choose a non-measurable $N \subseteq h(C)$. Then $h^{-1}(N) \subseteq C$ is a subset of a
null set, hence Lebesgue measurable. It is not a Borel set: otherwise $N$, its preimage under the continuous map $h^{-1}$, would be Borel, hence measurable.
In particular, a composite of a Lebesgue measurable function with a continuous one need not be Lebesgue measurable: $\mathbf{1}_{h^{-1}(N)} \circ h^{-1} = \mathbf{1}_N$.
\end{counterexample}

\section{Product measures and Fubini's theorem}\label{measure-theory:section-fubini}

\begin{definition}\label{measure-theory:definition-product-measure}
For measurable spaces $(X, \mathcal{A})$ and $(Y, \mathcal{B})$, the \emph{product $\sigma$-algebra} $\mathcal{A} \otimes \mathcal{B}$ on $X \times Y$ is generated by the \emph{rectangles} $A \times B$. For
$E \subseteq X \times Y$, the sections are $E_x = \{y : (x, y) \in E\}$ and $E^y = \{x : (x, y) \in E\}$.
\end{definition}

\begin{theorem}\label{measure-theory:theorem-product-measure}
Let $\mu$ and $\nu$ be $\sigma$-finite measures on $X$ and $Y$. For $E \in \mathcal{A} \otimes \mathcal{B}$, the sections are measurable, the functions $x \mapsto \nu(E_x)$ and $y \mapsto \mu(E^y)$ are measurable, and
\[
(\mu \times \nu)(E) = \int\nu(E_x)\,d\mu(x) = \int\mu(E^y)\,d\nu(y)
\]
defines the unique measure on $\mathcal{A} \otimes \mathcal{B}$ with $(\mu \times \nu)(A \times B) = \mu(A)\nu(B)$.
\end{theorem}

\begin{proof}
Assume first that $\mu$ and $\nu$ are finite. The sets $E$ for which all claims hold (including equality of the two integrals) contain the rectangles, which are closed under finite intersections, and
they are closed under complements and countable disjoint unions (by additivity and monotone convergence). Measurability of sections holds for all $E \in \mathcal{A} \otimes \mathcal{B}$ since the sets with
measurable sections form a $\sigma$-algebra. By Lemma~\ref{measure-theory:lemma-dynkin} the claims hold on $\mathcal{A} \otimes \mathcal{B}$, and countable additivity follows from monotone convergence.
For $\sigma$-finite measures write $X \times Y$ as an increasing union of $X_m \times Y_m$ with finite measures and pass to the limit. Uniqueness follows from Lemma~\ref{measure-theory:lemma-dynkin}.
\end{proof}

\begin{theorem}[Tonelli and Fubini]\label{measure-theory:theorem-fubini}
Let $\mu$, $\nu$ be $\sigma$-finite.
\begin{enumerate}
\item If $f \colon X \times Y \to [0, \infty]$ is measurable, then $x \mapsto \int f(x, y)\,d\nu(y)$ is measurable and
$\int_{X \times Y}f\,d(\mu \times \nu) = \int\left(\int f(x, y)\,d\nu(y)\right)d\mu(x) = \int\left(\int f(x, y)\,d\mu(x)\right)d\nu(y)$.
\item If $f \in L^1(\mu \times \nu)$, then $f(x, -) \in L^1(\nu)$ for a.e. $x$, the function $x \mapsto \int f(x, y)\,d\nu(y)$ (defined a.e.) is integrable, and the same equalities hold.
\end{enumerate}
\end{theorem}

\begin{proof}
(1) For indicator functions this is Theorem~\ref{measure-theory:theorem-product-measure}; it extends to simple functions by linearity and to all $f \geq 0$ by monotone convergence. (2) Apply (1) to $|f|$ to see that
$\int|f(x, y)|\,d\nu(y) < \infty$ for a.e. $x$, and then to the positive and negative parts of the real and imaginary parts.
\end{proof}

\begin{counterexample}\label{measure-theory:counterexample-fubini}
Integrability is needed in (2). On $\NN \times \NN$ with counting measures let $f(n, n) = 1$, $f(n, n+1) = -1$ and $f = 0$ otherwise. Then $\sum_m\sum_nf(n, m) = 1$, since only column $m = 0$ has nonzero
sum, whereas $\sum_n\sum_mf(n, m) = 0$.
\end{counterexample}

\begin{remark}\label{measure-theory:remark-lebesgue-product}
$\lambda^{m+n}$ and $\lambda^m \times \lambda^n$ agree on boxes, hence on Borel sets by Lemma~\ref{measure-theory:lemma-dynkin}, and $\lambda^{m+n}$ is the completion of $\lambda^m \times \lambda^n$. So
Fubini's theorem applies to Lebesgue measure on $\RR^{m+n} = \RR^m \times \RR^n$ for Borel functions, and for Lebesgue measurable functions with ``for almost every $x$'' qualifications.
\end{remark}

\begin{proposition}[Layer cake formula]\label{measure-theory:proposition-layer-cake}
Let $\mu$ be a $\sigma$-finite measure on $X$ and $f \colon X \to [0, \infty]$ measurable. Then
\[
\int_Xf\,d\mu = \int_0^\infty\mu(\{f > t\})\,dt.
\]
More generally, if $\psi \colon [0, \infty) \to [0, \infty)$ is Borel measurable and $\phi(s) = \int_0^s\psi(t)\,dt$ (with $\phi(\infty) = \int_0^\infty\psi$), then
$\int_X\phi \circ f\,d\mu = \int_0^\infty\psi(t)\,\mu(\{f > t\})\,dt$; in particular $\int_Xf^p\,d\mu = p\int_0^\infty t^{p-1}\mu(\{f > t\})\,dt$ for $p > 0$.
\end{proposition}

\begin{proof}
The set $E = \{(x, t) \in X \times [0, \infty) : t < f(x)\}$ is measurable for the product $\sigma$-algebra, because $t < f(x)$ holds exactly when $t < q < f(x)$ for some rational
$q$, so $E = \bigcup_{q \in \QQ_{>0}}\{f > q\} \times [0, q)$. Apply Tonelli's theorem (Theorem~\ref{measure-theory:theorem-fubini}(1)) to the non-negative measurable function
$(x, t) \mapsto \psi(t)\mathbf{1}_E(x, t)$ on $X \times [0, \infty)$ with $\mu \times \lambda$. Integrating first in $t$ gives $\int_0^{f(x)}\psi(t)\,dt = \phi(f(x))$, integrating first in
$x$ gives $\psi(t)\mu(\{f > t\})$. The first formula is the case $\psi = 1$, and the last one the case $\psi(t) = pt^{p-1}$.
\end{proof}

The name describes the picture: the region under the graph of $f$ is cut into horizontal layers $\{f > t\}$ of thickness $dt$. As an application we show that the maximal
function of an $L^p$ function is in $L^p$ for $p > 1$, although it is never in $L^1$ (Exercise~\ref{measure-theory:exercise-maximal-not-integrable}).

\begin{corollary}[Hardy--Littlewood maximal theorem]\label{measure-theory:corollary-maximal-lp}
For $1 < p \leq \infty$ and $f \in L^p(\RR^n)$, $\|Mf\|_p \leq C_{n,p}\|f\|_p$, with $C_{n,\infty} = 1$ and $C_{n,p}^p = 2^p3^np/(p - 1)$.
\end{corollary}

\begin{proof}
For $p = \infty$, every average of $|f|$ is at most $\|f\|_\infty$. Let $p < \infty$ and $t > 0$, and split $f = f_1 + f_2$ with $f_1 = f\mathbf{1}_{\{|f| > t/2\}}$. Then $f_1 \in L^1$,
since $|f_1| \leq (2/t)^{p-1}|f|^p$, and $Mf \leq Mf_1 + Mf_2 \leq Mf_1 + t/2$, so $\{Mf > t\} \subseteq \{Mf_1 > t/2\}$, and by
Theorem~\ref{measure-theory:theorem-maximal-inequality}, $\lambda(\{Mf > t\}) \leq \frac{2 \cdot 3^n}{t}\int_{\{|f| > t/2\}}|f|$. By
Proposition~\ref{measure-theory:proposition-layer-cake} and Tonelli's theorem,
\[
\int(Mf)^p = p\int_0^\infty t^{p-1}\lambda(\{Mf > t\})\,dt \leq 2 \cdot 3^np\int|f(x)|\int_0^{2|f(x)|}t^{p-2}\,dt\,d\lambda(x) = \frac{2 \cdot 3^np}{p - 1}\int|f|(2|f|)^{p-1},
\]
which is $\frac{2^p3^np}{p - 1}\|f\|_p^p$.
\end{proof}

\section{Change of variables}\label{measure-theory:section-change-of-variables}

\begin{lemma}\label{measure-theory:lemma-linear-change}
For $T \in \GL_n(\RR)$ and Borel $E \subseteq \RR^n$, $\lambda(T(E)) = |\det T|\,\lambda(E)$.
\end{lemma}

\begin{proof}
Both sides are multiplicative in $T$, and every invertible matrix is a product of \emph{elementary} matrices (by Gaussian elimination): permutations of coordinates, scalings of one coordinate by
$c \neq 0$, and shears $x_i \mapsto x_i + cx_j$ ($i \neq j$). Permutations and scalings map boxes to boxes, multiplying volumes by $|\det T|$, so they multiply outer measure by $|\det T|$. A shear preserves each
line parallel to the $i$-th axis and acts on it by a translation; by Theorem~\ref{measure-theory:theorem-fubini} and Remark~\ref{measure-theory:remark-lebesgue-product}, integrating the one-dimensional
measures of these sections shows $\lambda(T(E)) = \lambda(E)$, and $\det T = 1$.
\end{proof}

\begin{theorem}[Change of variables]\label{measure-theory:theorem-change-of-variables}
Let $U, V \subseteq \RR^n$ be open and $\phi \colon U \to V$ a $C^1$-diffeomorphism. For every measurable $f \colon V \to [0, \infty]$, and every $f \in L^1(V)$,
\[
\int_Vf\,d\lambda = \int_Uf(\phi(x))\,|\det D\phi(x)|\,d\lambda(x).
\]
\end{theorem}

\begin{proof}
We use the norm $\|x\|_\infty$ and the corresponding operator norm; the mean value inequality (Proposition~\ref{real-analysis:proposition-derivative-basic}(4)) holds for it with the same proof. Cubes are balls
$Q = \{x : \|x - a\|_\infty \leq h\}$.

\emph{Step 1.} Let $Q \subseteq U$ be a cube with centre $a$ and $T \in \GL_n(\RR)$, and let $M = \sup_{y \in Q}\|T^{-1}D\phi(y)\|$. By the mean value inequality applied to $T^{-1}\phi$,
$\phi(Q) \subseteq \phi(a) + T(\{\|z\|_\infty \leq Mh\})$, so $\lambda(\phi(Q)) \leq |\det T|M^n\lambda(Q)$ by Lemma~\ref{measure-theory:lemma-linear-change}.

\emph{Step 2: $\lambda(\phi(Q)) \leq \int_Q|\det D\phi|$ for cubes $Q \subseteq U$.} Let $\varepsilon > 0$. Since $D\phi$ is uniformly continuous on $Q$ and $\|D\phi^{-1}\|$ is bounded there, we can divide $Q$ into
subcubes $Q_j$ with centres $x_j$ such that $\|D\phi(x_j)^{-1}D\phi(y)\| \leq 1 + \varepsilon$ for $y \in Q_j$, and $|\det D\phi|$ varies by less than $\varepsilon$ on each $Q_j$. Step 1 with $T = D\phi(x_j)$ gives
$\lambda(\phi(Q)) \leq \sum_j\lambda(\phi(Q_j)) \leq (1 + \varepsilon)^n\sum_j|\det D\phi(x_j)|\lambda(Q_j) \leq (1 + \varepsilon)^n(\int_Q|\det D\phi| + \varepsilon\lambda(Q))$.

\emph{Step 3: $\lambda(\phi(E)) \leq \int_E|\det D\phi|$ for measurable $E \subseteq U$.} An open $W \subseteq U$ is a countable union of cubes with disjoint interiors (dyadic cubes), whose boundaries are null, so the
inequality holds for $W$. Let $E \subseteq K$ with $K \subseteq U$ compact, and $C = \sup$ of $|\det D\phi|$ on a compact neighbourhood $N \subseteq U$ of $K$. By regularity choose open $W$ with $E \subseteq W \subseteq N$ and
$\lambda(W \setminus E) < \delta$; then $\lambda(\phi(E)) \leq \lambda(\phi(W)) \leq \int_E|\det D\phi| + C\delta$. This shows that $\phi$ maps null sets to null sets, hence Lebesgue measurable sets to
Lebesgue measurable sets (a measurable set is a Borel set up to a null set), and gives the inequality for $E$ in a compact subset; exhaust $U$ by compact sets for general $E$.

\emph{Step 4.} By Step 3, $\int_V\mathbf{1}_{\phi(E)} \leq \int_U\mathbf{1}_{\phi(E)} \circ \phi\,|\det D\phi|$; by linearity and monotone convergence, $\int_Vf \leq \int_U(f \circ \phi)|\det D\phi|$ for all measurable
$f \geq 0$. Applying this to $\phi^{-1}$ and the function $(f \circ \phi)|\det D\phi|$ on $U$, and using $D\phi(\phi^{-1}(y))D(\phi^{-1})(y) = 1$, gives the reverse inequality. The $L^1$ case follows by linearity.
\end{proof}

\begin{example}\label{measure-theory:example-gaussian}
With polar coordinates $(r, \theta) \mapsto (r\cos\theta, r\sin\theta)$, a diffeomorphism $(0, \infty) \times (0, 2\pi) \to \RR^2$ minus a closed half-line (a null set), with Jacobian determinant $r$,
Tonelli's theorem gives $(\int_\RR e^{-x^2}dx)^2 = \int_{\RR^2}e^{-\|x\|^2} = \int_0^{2\pi}\int_0^\infty e^{-r^2}r\,dr\,d\theta = \pi$. So $\int_\RR e^{-x^2}\,dx = \sqrt{\pi}$.
\end{example}

Fubini's theorem reduces integration by parts in several variables to one variable. When one of the functions has compact support there is no boundary term.

\begin{proposition}[Integration by parts in $\RR^n$]\label{measure-theory:proposition-integration-by-parts}
Let $f, g \colon \RR^n \to \CC$ be continuously differentiable, and suppose that one of them has compact support. Then for each $i$
\[
\int_{\RR^n}(\partial_if)\,g\,d\lambda = -\int_{\RR^n}f\,\partial_ig\,d\lambda .
\]
Consequently, for a continuously differentiable vector field $F$ and a continuously differentiable function $g$, one of them compactly supported,
$\int\operatorname{div}F\,g\,d\lambda = -\int\langle F, \nabla g\rangle\,d\lambda$, where $\operatorname{div}F = \sum_i\partial_iF^i$ and $\nabla g = (\partial_1g, \ldots, \partial_ng)$; and for $f$ of class $C^2$,
$\int(\sum_i\partial_i^2f)\,g\,d\lambda = -\int\langle\nabla f, \nabla g\rangle\,d\lambda$.
\end{proposition}

\begin{proof}
The product $h = fg$ is continuously differentiable with compact support, contained in a cube $[-R, R]^n$, and $\partial_ih = (\partial_if)g + f\,\partial_ig$ by the product rule; all these functions are
continuous with compact support, hence integrable. By Fubini's theorem (Theorem~\ref{measure-theory:theorem-fubini}) we may integrate first in the variable $x_i$, the others being fixed. For fixed
values of the other variables, $t \mapsto h(\ldots, t, \ldots)$ is continuously differentiable and vanishes for $|t| \geq R$, so by the fundamental theorem of calculus
(Theorem~\ref{real-analysis:theorem-ftc}) $\int_{-R}^R\partial_ih\,dx_i = h(\ldots, R, \ldots) - h(\ldots, -R, \ldots) = 0$. Hence $\int\partial_ih\,d\lambda = 0$, which is the formula. For the vector field,
apply it to $F^i$ and $g$ and sum over $i$; for the last formula take $F = \nabla f$.
\end{proof}

The formula says that the operator $\partial_i$ is \emph{anti-self-adjoint} for the $L^2$ inner product on compactly supported functions: $\langle\partial_if, \bar g\rangle = -\langle f, \partial_i\bar g\rangle$. It is the
reason why the derivative of a distribution is defined by $(\partial_iT)(\varphi) = -T(\partial_i\varphi)$ (Chapter~\ref{distributions}), and, applied twice, it shows that $-\sum_i\partial_i^2$ is a
non-negative operator: $\int(-\sum_i\partial_i^2f)\bar f\,d\lambda = \int|\nabla f|^2\,d\lambda \geq 0$ for compactly supported $f$ of class $C^2$.

\section{The Radon--Nikodym and Riesz representation theorems}\label{measure-theory:section-radon-nikodym}

\begin{lemma}\label{measure-theory:lemma-hilbert-riesz}
Let $H$ be a complex inner product space that is complete (a \emph{Hilbert space}), for instance $L^2(\mu)$. Every continuous linear functional $\Lambda \colon H \to \CC$ is of the form $\Lambda(f) = \langle f, g \rangle$
for a unique $g \in H$.
\end{lemma}

\begin{proof}
If $\Lambda = 0$ take $g = 0$. Otherwise $N = \ker\Lambda$ is a closed proper subspace. Choose $u \notin N$ and let $d = \inf_{v \in N}\|u - v\| > 0$. For $v_k \in N$ with $\|u - v_k\| \to d$, the parallelogram law
$\|v_k - v_l\|^2 = 2\|u - v_k\|^2 + 2\|u - v_l\|^2 - 4\|u - \frac{v_k + v_l}{2}\|^2$ shows that $(v_k)$ is Cauchy; its limit $v$ satisfies $\|u - v\| = d$. For $w \in N$ and $t \in \CC$,
$\|u - v - tw\|^2 \geq d^2$ forces $\langle u - v, w \rangle = 0$. So $z = u - v \neq 0$ is orthogonal to $N$. For $f \in H$, $\Lambda(f)z - \Lambda(z)f \in N$, so $\Lambda(f)\|z\|^2 = \Lambda(z)\langle f, z \rangle$, and
$g = \overline{\Lambda(z)}z/\|z\|^2$ works. Uniqueness: if $\langle f, g - g' \rangle = 0$ for all $f$, take $f = g - g'$.
\end{proof}

\begin{theorem}[Radon--Nikodym]\label{measure-theory:theorem-radon-nikodym}
Let $\mu$ and $\nu$ be $\sigma$-finite measures on $(X, \mathcal{A})$ with $\nu \ll \mu$, meaning $\mu(E) = 0 \Rightarrow \nu(E) = 0$. There is a measurable $h \geq 0$, unique a.e., with
$\nu(E) = \int_Eh\,d\mu$ for all $E$.
\end{theorem}

\begin{proof}
Assume $\mu$ and $\nu$ finite; the $\sigma$-finite case follows by decomposing $X$. Let $\rho = \mu + \nu$. On $L^2(\rho)$, $f \mapsto \int f\,d\nu$ is continuous, since
$|\int f\,d\nu| \leq \int|f|\,d\rho \leq \rho(X)^{1/2}\|f\|_2$. By Lemma~\ref{measure-theory:lemma-hilbert-riesz} there is $g$ with $\int f\,d\nu = \int fg\,d\rho$ for all $f \in L^2(\rho)$, and $g$ may be taken real. Taking
$f = \mathbf{1}_E$ gives $\nu(E) = \int_Eg\,d\rho$, so $g \geq 0$ $\rho$-a.e.; and $\int_E(1 - g)\,d\nu = \int_Eg\,d\mu$ (for bounded $f$, $\int f(1 - g)\,d\nu = \int fg\,d\mu$). For $E = \{g \geq 1\}$ the left side
is $\leq 0$ and the right side is $\geq \mu(E)$, so $\mu(E) = 0$ and $\nu(E) = 0$. Applying the identity to $f = \mathbf{1}_E(1 + g + \cdots + g^k)$ gives $\int_E(1 - g^{k+1})\,d\nu = \int_E(g + \cdots + g^{k+1})\,d\mu$. As
$k \to \infty$, monotone convergence gives $\nu(E) = \int_Eh\,d\mu$ with $h = g/(1 - g)$. If $h'$ also works, $\int_E(h - h')\,d\mu = 0$ for $E = \{h > h'\}$ and $\{h < h'\}$ (intersected with sets of finite measure), so
$h = h'$ a.e.
\end{proof}

\begin{theorem}[Riesz--Markov]\label{measure-theory:theorem-riesz-markov}
Let $X$ be a locally compact Hausdorff space and $\Lambda \colon C_c(X, \RR) \to \RR$ a linear map with $\Lambda(f) \geq 0$ for $f \geq 0$. There is a unique regular Borel measure $\mu$, finite on compact sets, with
$\Lambda(f) = \int f\,d\mu$ for all $f \in C_c(X, \RR)$.
\end{theorem}

This is \cite[Theorem 2.14]{rudin-rca}. For $X = \RR^n$ and $\Lambda$ the Riemann integral of continuous functions with compact support, $\mu$ is Lebesgue measure restricted to Borel sets. On a smooth manifold,
a volume form defines such a functional, and hence a measure (Chapter~\ref{differential-forms}).

\section{Absolutely continuous functions}\label{measure-theory:section-absolutely-continuous}

By Corollary~\ref{measure-theory:corollary-ftc-lebesgue}, an indefinite integral $F(x) = F(a) + \int_a^xf$ of an integrable function satisfies $F' = f$ almost everywhere.
Conversely, when is a function the integral of its derivative? Continuity and differentiability almost everywhere are not enough, by the Cantor function
(Counterexample~\ref{measure-theory:counterexample-cantor-function}). The right condition is absolute continuity.

\begin{definition}\label{measure-theory:definition-absolutely-continuous}
A function $F \colon [a, b] \to \CC$ is \emph{absolutely continuous} if for every $\varepsilon > 0$ there is $\delta > 0$ such that $\sum_k|F(b_k) - F(a_k)| < \varepsilon$ for every finite
family of pairwise disjoint intervals $(a_k, b_k) \subseteq [a, b]$ with $\sum_k(b_k - a_k) < \delta$. For real $F$, the \emph{total variation} of $F$ on $[a, x]$ is
$V_F(x) = \sup\sum_i|F(t_i) - F(t_{i-1})|$, the supremum over the partitions $a = t_0 < \cdots < t_m = x$.
\end{definition}

\begin{lemma}\label{measure-theory:lemma-absolutely-continuous}
\begin{enumerate}
\item For $f \in L^1([a, b])$, $F(x) = \int_{[a,x]}f\,d\lambda$ is absolutely continuous; more precisely, for every $\varepsilon > 0$ there is $\delta > 0$ with $\int_E|f| < \varepsilon$
whenever $\lambda(E) < \delta$.
\item If $F$ is real and absolutely continuous, then $V_F$ is finite, nondecreasing and absolutely continuous, and so is $V_F - F$. So $F = V_F - (V_F - F)$ is a difference of
nondecreasing absolutely continuous functions.
\end{enumerate}
\end{lemma}

\begin{proof}
(1) By monotone convergence $\int_{\{|f| > N\}}|f| \to 0$ as $N \to \infty$; choose $N$ with $\int_{\{|f| > N\}}|f| < \varepsilon/2$ and $\delta = \varepsilon/(2N)$. Then
$\int_E|f| \leq \int_{\{|f| > N\}}|f| + N\lambda(E) < \varepsilon$. For disjoint intervals of total length $< \delta$, $\sum_k|F(b_k) - F(a_k)| \leq \int_{\bigcup_k(a_k, b_k)}|f| < \varepsilon$.

(2) Take $\delta$ for $\varepsilon = 1$ and cut $[a, b]$ into $N$ intervals of length $< \delta$; refining a partition by these cut points does not decrease the sum, and the sum
over each of the $N$ intervals is $< 1$, so $V_F(b) \leq N$. For $x < y$, $V_F(y) - V_F(x)$ is the total variation of $F$ on $[x, y]$ (a partition of $[a, y]$ refined by $x$ splits
into partitions of $[a, x]$ and $[x, y]$), so $V_F$ is nondecreasing and $V_F(y) - V_F(x) \geq |F(y) - F(x)| \geq F(y) - F(x)$, which makes $V_F - F$ nondecreasing. If
$(a_k, b_k)$ are disjoint with $\sum_k(b_k - a_k) < \delta$ for the $\delta$ of $\varepsilon$, then any partitions of the intervals $[a_k, b_k]$ together give disjoint intervals of total
length $< \delta$, so $\sum_k(V_F(b_k) - V_F(a_k)) \leq \varepsilon$: $V_F$ is absolutely continuous, and so is $V_F - F$.
\end{proof}

\begin{lemma}\label{measure-theory:lemma-stieltjes}
Let $F \colon [a, b] \to \RR$ be continuous and nondecreasing. There is a finite Borel measure $\mu_F$ on $[a, b]$ with $\mu_F([a, x]) = F(x) - F(a)$ for all $x$; so
$\mu_F((x, y]) = F(y) - F(x)$ and points have measure $0$. If $F$ is absolutely continuous, then $\mu_F$ is absolutely continuous with respect to $\lambda$.
\end{lemma}

\begin{proof}
Let $c = F(a)$ and $d = F(b)$. For $t \in [c, d]$ the set $\{x : F(x) = t\}$ is nonempty by the intermediate value theorem and closed, so it has a least element $G(t)$. The map
$G \colon [c, d] \to [a, b]$ is nondecreasing, hence Borel measurable (preimages of intervals are intervals), and $G(t) \leq x$ if and only if $t \leq F(x)$: if $t \leq F(x)$ some
$y \leq x$ has $F(y) = t$, and if $G(t) \leq x$ then $t = F(G(t)) \leq F(x)$. Let $\mu_F(E) = \lambda(\{t \in [c, d] : G(t) \in E\})$ for Borel $E \subseteq [a, b]$, the image of
Lebesgue measure under $G$. Then $\mu_F([a, x]) = \lambda([c, F(x)]) = F(x) - c$.

Let $F$ be absolutely continuous and $E$ a Borel set with $\lambda(E) = 0$. Given $\varepsilon > 0$, take $\delta$ from the definition and an open $U \supseteq E$ with $\lambda(U) < \delta$
(Theorem~\ref{measure-theory:theorem-lebesgue-measure}). The open set $U \cap (a, b)$ is the disjoint union of countably many open intervals $(a_k, b_k)$, its connected components,
with $\sum_k(b_k - a_k) < \delta$. For every $N$, $\sum_{k \leq N}\mu_F((a_k, b_k)) = \sum_{k \leq N}(F(b_k) - F(a_k)) < \varepsilon$, so $\mu_F(U \cap (a, b)) \leq \varepsilon$ and
$\mu_F(E) \leq \varepsilon$, the endpoints $a$, $b$ having measure $0$. As $\varepsilon$ is arbitrary, $\mu_F(E) = 0$.
\end{proof}

\begin{theorem}[Fundamental theorem of calculus for absolutely continuous functions]\label{measure-theory:theorem-absolutely-continuous}
Let $F \colon [a, b] \to \CC$ be absolutely continuous. Then $F$ is differentiable almost everywhere, $F' \in L^1([a, b])$, and
\[
F(x) = F(a) + \int_a^xF'(t)\,dt \qquad \text{for all } x \in [a, b].
\]
Conversely, every function of the form $F(a) + \int_a^xf$ with $f \in L^1([a, b])$ is absolutely continuous. So the absolutely continuous functions are exactly the indefinite
Lebesgue integrals.
\end{theorem}

\begin{proof}
The converse is Lemma~\ref{measure-theory:lemma-absolutely-continuous}(1). For the direct statement, treat real and imaginary parts separately, and by
Lemma~\ref{measure-theory:lemma-absolutely-continuous}(2) assume that $F$ is nondecreasing, the statement being linear in $F$. By
Lemma~\ref{measure-theory:lemma-stieltjes}, $\mu_F$ is a finite measure with $\mu_F \ll \lambda$ on the Borel sets of $[a, b]$, so by the Radon--Nikodym theorem
(Theorem~\ref{measure-theory:theorem-radon-nikodym}) there is an integrable $h \geq 0$ with $\mu_F(E) = \int_Eh\,d\lambda$. Hence $F(x) - F(a) = \mu_F([a, x]) = \int_a^xh$, and by
Corollary~\ref{measure-theory:corollary-ftc-lebesgue}, $F$ is differentiable almost everywhere with $F' = h$.
\end{proof}

\begin{example}\label{measure-theory:example-absolutely-continuous}
A Lipschitz function, $|F(x) - F(y)| \leq L|x - y|$, is absolutely continuous (take $\delta = \varepsilon/L$), so it is differentiable almost everywhere and is the integral of its
derivative; for instance $|x| = \int_0^x\operatorname{sgn}$. The Cantor function is continuous and nondecreasing but not absolutely continuous: it maps the Cantor set, of measure
$0$, onto $[0, 1]$, whereas the proof of Lemma~\ref{measure-theory:lemma-stieltjes} shows that an absolutely continuous nondecreasing function gives measure $0$ to null sets, and
$\mu_F(C) = F(1) - F(0) = 1$ for the Cantor function because $F$ is constant on the complement of $C$. The function $x\sin(1/x)$ (and $0$ at $0$) on $[0, 1]$ is continuous but not
of bounded variation, hence not absolutely continuous.
\end{example}


\section{Hausdorff measures}\label{measure-theory:section-hausdorff}

Lebesgue measure $\lambda^n$ measures $n$-dimensional volume in $\RR^n$, but it does not see lower-dimensional sets: a curve in the plane, or a surface in space, has measure $0$. How long is
a curve, and what is the area of a surface, without choosing a parametrisation? And how big is the Cantor set, which has measure $0$ and is uncountable? Hausdorff measures answer these
questions at once: for every $s \geq 0$ there is an $s$-dimensional measure $\mathcal{H}^s$ on any metric space, defined by covering a set with small sets $C$ and adding up
$(\operatorname{diam}C)^s$. This and the next two sections are an introduction to \emph{geometric measure theory}: Hausdorff measure and dimension, differentiability of Lipschitz maps,
and the area formula, which computes $\mathcal{H}^k$ of parametrised $k$-dimensional sets. We follow \cite[Chapters 2--3]{evans-gariepy}; the comprehensive reference is
\cite{federer-gmt}.

Throughout, $X$ is a metric space, $\operatorname{diam}C = \sup\{d(x, y) : x, y \in C\}$ for $C \neq \emptyset$, and $\operatorname{diam}\emptyset = 0$. For $s \geq 0$ let
$\alpha(s) = \pi^{s/2}/(2^s\Gamma(\frac s2 + 1))$, with $\Gamma$ as in Exercise~\ref{measure-theory:exercise-ball-volume}; for an integer $n$, $\alpha(n)$ is the volume of a ball of
diameter $1$ in $\RR^n$. The constant only serves to make $\mathcal{H}^n$ equal to Lebesgue measure on $\RR^n$ (Theorem~\ref{measure-theory:theorem-hausdorff-lebesgue}); in the sums below
the empty set contributes $0$ and every nonempty set $C$ contributes $\alpha(s)(\operatorname{diam}C)^s$, with the convention $(\operatorname{diam}C)^0 = 1$, also when $C$ is a
single point.

\begin{definition}\label{measure-theory:definition-hausdorff}
For $A \subseteq X$, $s \geq 0$ and $\delta > 0$ let
\[
\mathcal{H}^s_\delta(A) = \inf\Bigl\{\sum_k\alpha(s)(\operatorname{diam}C_k)^s : A \subseteq \bigcup_kC_k,\ \operatorname{diam}C_k \leq \delta\Bigr\},
\]
the infimum over countable covers by sets of diameter at most $\delta$. As $\delta$ decreases, fewer covers are allowed and $\mathcal{H}^s_\delta(A)$ increases; the
\emph{$s$-dimensional Hausdorff (outer) measure} is $\mathcal{H}^s(A) = \lim_{\delta \to 0}\mathcal{H}^s_\delta(A) = \sup_{\delta > 0}\mathcal{H}^s_\delta(A)$.
\end{definition}

The restriction $\operatorname{diam}C_k \leq \delta$ is essential: with arbitrary covers, a long segment could be covered by a single set and would get the same size as a short one, while
small sets are forced to follow the shape of $A$. Replacing each $C_k$ by its closure does not change its diameter, so closed covers give the same infimum. And $\mathcal{H}^s(A)$ only
depends on the metric space $A$: a cover of $A$ by subsets of $X$ gives the cover by the sets $C_k \cap A$, of smaller diameters. In particular an isometry $A \to B$ between subsets of two
metric spaces preserves $\mathcal{H}^s$.

\begin{lemma}\label{measure-theory:lemma-metric-outer-measure}
Let $\mu^*$ be an outer measure on $X$ which is \emph{metric}: $\mu^*(A \cup B) = \mu^*(A) + \mu^*(B)$ whenever $d(A, B) = \inf\{d(a, b) : a \in A, b \in B\} > 0$. Then every Borel set is
$\mu^*$-measurable.
\end{lemma}

\begin{proof}
The measurable sets form a $\sigma$-algebra (Theorem~\ref{measure-theory:theorem-caratheodory}), so it suffices to show that a closed set $F$ is measurable, that is,
$\mu^*(E \cap F) + \mu^*(E \setminus F) \leq \mu^*(E)$ for every $E$ with $\mu^*(E) < \infty$. Let $E_j = \{x \in E : d(x, F) \geq 1/j\}$. Since $d(E \cap F, E_j) \geq 1/j$,
$\mu^*(E \cap F) + \mu^*(E_j) = \mu^*((E \cap F) \cup E_j) \leq \mu^*(E)$, and it remains to show $\mu^*(E_j) \to \mu^*(E \setminus F)$. Let
$R_j = \{x \in E : \frac1{j+1} \leq d(x, F) < \frac1j\}$. As $F$ is closed, $E \setminus F = E_j \cup \bigcup_{i \geq j}R_i$. If $i \geq l + 2$, then $d(R_i, R_l) > 0$, so by induction
$\sum_{i \leq N}\mu^*(R_{2i}) = \mu^*(\bigcup_{i \leq N}R_{2i}) \leq \mu^*(E)$, and similarly for the odd indices. Hence $\sum_i\mu^*(R_i) < \infty$, and
$\mu^*(E \setminus F) \leq \mu^*(E_j) + \sum_{i \geq j}\mu^*(R_i)$, where the tail tends to $0$. With $\mu^*(E_j) \leq \mu^*(E \setminus F)$ this gives the limit.
\end{proof}

\begin{proposition}\label{measure-theory:proposition-hausdorff}
\begin{enumerate}
\item $\mathcal{H}^s$ is a metric outer measure, so its restriction to the Borel sets is a measure.
\item $\mathcal{H}^0$ is the counting measure.
\item If $f \colon A \to Y$ is Lipschitz with constant $L$, then $\mathcal{H}^s(f(A)) \leq L^s\mathcal{H}^s(A)$. For $X = \RR^n$, $\mathcal{H}^s$ is translation invariant and
$\mathcal{H}^s(rA) = r^s\mathcal{H}^s(A)$ for $r > 0$.
\item If $\mathcal{H}^s(A) < \infty$ and $t > s$, then $\mathcal{H}^t(A) = 0$. If $\mathcal{H}^t(A) > 0$ and $s < t$, then $\mathcal{H}^s(A) = \infty$.
\end{enumerate}
\end{proposition}

\begin{proof}
(1) Each $\mathcal{H}^s_\delta$ is an outer measure: it is monotone, and given covers of the $A_k$ with sums at most $\mathcal{H}^s_\delta(A_k) + \varepsilon2^{-k}$, their union covers
$\bigcup_kA_k$. Then $\mathcal{H}^s(\bigcup_kA_k) = \sup_\delta\mathcal{H}^s_\delta(\bigcup_kA_k) \leq \sup_\delta\sum_k\mathcal{H}^s_\delta(A_k) \leq \sum_k\mathcal{H}^s(A_k)$. If
$d(A, B) > 0$ and $\delta < d(A, B)$, no set of diameter at most $\delta$ meets both $A$ and $B$, so a $\delta$-cover of $A \cup B$ splits into covers of $A$ and of $B$, and
$\mathcal{H}^s_\delta(A \cup B) \geq \mathcal{H}^s_\delta(A) + \mathcal{H}^s_\delta(B)$; let $\delta \to 0$. Apply Lemma~\ref{measure-theory:lemma-metric-outer-measure}.

(2) $\alpha(0) = 1$. A set of $N$ points needs at least $N$ nonempty sets of diameter less than the smallest distance between the points, and the $N$ singletons suffice; so
$\mathcal{H}^0_\delta$ of it is $N$ for small $\delta$. An infinite set contains finite sets of every size.

(3) If $(C_k)$ is a $\delta$-cover of $A$, then $(f(C_k \cap A))$ is an $L\delta$-cover of $f(A)$ with $\operatorname{diam}f(C_k \cap A) \leq L\operatorname{diam}C_k$; so
$\mathcal{H}^s_{L\delta}(f(A)) \leq L^s\mathcal{H}^s_\delta(A)$. Translations are isometries, and $x \mapsto rx$ multiplies diameters by $r$, with inverse $x \mapsto x/r$.

(4) For a $\delta$-cover, $\sum_k\alpha(t)(\operatorname{diam}C_k)^t \leq \frac{\alpha(t)}{\alpha(s)}\delta^{t-s}\sum_k\alpha(s)(\operatorname{diam}C_k)^s$, so
$\mathcal{H}^t_\delta(A) \leq \frac{\alpha(t)}{\alpha(s)}\delta^{t-s}\mathcal{H}^s_\delta(A) \leq \frac{\alpha(t)}{\alpha(s)}\delta^{t-s}\mathcal{H}^s(A) \to 0$. The second statement is the
contrapositive.
\end{proof}

By (4), for each $A$ there is a unique $s_0 \in [0, \infty]$ with $\mathcal{H}^s(A) = \infty$ for $s < s_0$ and $\mathcal{H}^s(A) = 0$ for $s > s_0$; it is the \emph{Hausdorff dimension}
$\dim_{\mathcal{H}}A = \inf\{s : \mathcal{H}^s(A) = 0\}$. At $s = s_0$ the measure can be anything in $[0, \infty]$. By (3), Lipschitz maps do not increase Hausdorff dimension, and
bi-Lipschitz maps preserve it.

\begin{theorem}\label{measure-theory:theorem-hausdorff-lebesgue}
On the Borel sets of $\RR^n$, $\mathcal{H}^n = \lambda^n$. Consequently the subsets of $\RR^n$ have Hausdorff dimension at most $n$, and nonempty open subsets have dimension $n$.
\end{theorem}

\begin{proof}
We show that $\mathcal{H}^n = c\lambda^n$ for a constant $0 < c < \infty$; that $c = 1$ is the content of the \emph{isodiametric inequality} $\lambda^n(C) \leq \alpha(n)(\operatorname{diam}C)^n$,
which says that among sets of given diameter the balls have the largest volume, together with Vitali's covering theorem, and we refer to \cite[Section 2.2]{evans-gariepy} for these two steps.

\emph{The constant.} Let $Q = [0, 1)^n$. Covering $Q$ by $m^n$ cubes of side $1/m$ and diameter $\sqrt n/m$ gives $\mathcal{H}^n_{\sqrt n/m}(Q) \leq \alpha(n)n^{n/2}$, so
$\mathcal{H}^n(Q) < \infty$. Conversely, a set $C$ of diameter $d$ lies in a closed box with sides $d$ (each coordinate varies by at most $d$ on $C$), so $\lambda^*(C) \leq d^n$; for any cover
of $Q$, $\sum_k\alpha(n)(\operatorname{diam}C_k)^n \geq \alpha(n)\sum_k\lambda^*(C_k) \geq \alpha(n)$, so $c = \mathcal{H}^n(Q) \geq \alpha(n) > 0$.

\emph{Comparison.} $\mathcal{H}^n$ is a translation invariant measure on the Borel sets (Proposition~\ref{measure-theory:proposition-hausdorff}). For $j \geq 0$, $Q$ is the disjoint union of $2^{jn}$
translates of $[0, 2^{-j})^n$, so every half-open dyadic cube $\prod_i[a_i2^{-j}, (a_i + 1)2^{-j})$ ($a_i \in \ZZ$, $j \in \ZZ$) has $\mathcal{H}^n$-measure $c2^{-jn} = c\lambda^n$ of it; for
$j < 0$ use that such a cube is a disjoint union of $2^{-jn}$ unit cubes. Fix $m \geq 0$ and $P = [-2^m, 2^m)^n$. The dyadic cubes contained in $P$, together with $P$, form a family closed under
finite intersections, on which $\mathcal{H}^n$ and $c\lambda^n$ agree and are finite. By Lemma~\ref{measure-theory:lemma-dynkin} the two measures agree on the $\sigma$-algebra it generates, which
contains every open subset of the interior of $P$ (a countable union of dyadic cubes), hence every Borel subset of the interior. The boundary of $P$ is a finite union of pieces of hyperplanes,
which are null for both measures: a piece $[0, 1]^{n-1} \times \{0\}$ is covered by $j^{n-1}$ cubes of diameter $\sqrt n/j$, with $\sum\alpha(n)(\operatorname{diam})^n \to 0$ as $j \to \infty$. So
the two measures agree on the Borel subsets of $P$, and letting $m \to \infty$, on all Borel sets.

The last statement follows from Proposition~\ref{measure-theory:proposition-hausdorff}(4): $\mathcal{H}^n$ is finite on bounded sets and $\RR^n$ is a countable union of them, so
$\mathcal{H}^t = 0$ for $t > n$, while a nonempty open set has positive $\mathcal{H}^n$-measure.
\end{proof}

\begin{example}\label{measure-theory:example-cantor-dimension}
\emph{The Cantor set has Hausdorff dimension $s = \log 2/\log 3$}, and $\alpha(s)/2 \leq \mathcal{H}^s(C) \leq \alpha(s)$; note $3^s = 2$. Let $C \subseteq [0, 1]$ be as in
Counterexample~\ref{measure-theory:counterexample-cantor-function}. For $k \geq 0$ and digits $a_1, \ldots, a_k \in \{0, 2\}$ let $I_{a_1\cdots a_k}$ be the closed interval with left end
$\sum_{j \leq k}a_j3^{-j}$ and length $3^{-k}$; there are $2^k$ of them, and $C$ is contained in their union $C_k$, since the digits beyond the $k$-th contribute at most $3^{-k}$.
\begin{enumerate}
\item \emph{Upper bound.} The $2^k$ intervals form a $3^{-k}$-cover of $C$ with $\sum\alpha(s)3^{-ks} = \alpha(s)2^k3^{-ks} = \alpha(s)$. So $\mathcal{H}^s(C) \leq \alpha(s)$.
\item \emph{Separation.} Two of these intervals of the same generation $k$ are at distance at least $3^{-k}$: if their digit strings first differ at position $j \leq k$, their left ends differ by
at least $2 \cdot 3^{-j} - \sum_{j < i \leq k}2 \cdot 3^{-i} = 3^{-j} + 3^{-k} \geq 2 \cdot 3^{-k}$. So a set of diameter less than $3^{-k}$ meets at most one of them.
\item \emph{The Cantor measure.} Let $\mu = \mu_F$ be the measure of Lemma~\ref{measure-theory:lemma-stieltjes} for the Cantor function $F$. Since $F$ is constant on each interval of
$[0, 1] \setminus C$, $\mu([0, 1] \setminus C) = 0$ and $\mu(C) = F(1) - F(0) = 1$. On the generation-$k$ interval $I$ with left end $a = \sum_{j \leq k}a_j3^{-j}$ and right end
$a + \sum_{j > k}2 \cdot 3^{-j}$, $\mu(I) = F(a + 3^{-k}) - F(a) = \sum_{j > k}2^{-j} = 2^{-k}$, by the formula for $c$ on $C$.
\item \emph{Lower bound.} Let $U$ be a closed set with $0 < \operatorname{diam}U = d < 1$, and $k$ with $3^{-k-1} \leq d < 3^{-k}$. By (2), $U$ meets at most one generation-$k$ interval $I$,
and $\mu(U) \leq \mu(I) + \mu([0, 1] \setminus C) = 2^{-k} = 3^s3^{-(k+1)s} \leq 2d^s$; for a single point, $\mu(U) = 0$. For a cover of $C$ by closed sets $U_i$ of diameter less than $1$,
$1 = \mu(C) \leq \sum_i\mu(U_i \cap [0, 1]) \leq 2\sum_i(\operatorname{diam}U_i)^s$. Hence $\mathcal{H}^s_\delta(C) \geq \alpha(s)/2$ for $\delta < 1$.
\end{enumerate}
The argument in (4), bounding the mass that a measure on $A$ gives to small sets, is the \emph{mass distribution principle}, the standard way to bound Hausdorff measures from below. So the Cantor
set is too small to have positive length but too large to have zero $s$-dimensional measure for any $s < \log 2/\log 3 \approx 0.63$.
\end{example}

\section{Lipschitz maps and Rademacher's theorem}\label{measure-theory:section-rademacher}

A Lipschitz function of one variable is absolutely continuous, hence differentiable almost everywhere (Example~\ref{measure-theory:example-absolutely-continuous}). Rademacher's theorem extends
this to several variables. It is the reason why Lipschitz maps are the natural class of maps in geometric measure theory: they are just regular enough to have a derivative almost everywhere,
and they are preserved by the operations of the theory (distance functions, compositions, restrictions, inverses of bi-Lipschitz maps), which smooth maps are not.

\begin{lemma}\label{measure-theory:lemma-vanishing-integrals}
If $g \in L^1_{\mathrm{loc}}(\RR^n)$ and $\int g\varphi\,d\lambda = 0$ for all compactly supported functions $\varphi$ of class $C^1$, then $g = 0$ almost everywhere.
\end{lemma}

\begin{proof}
Let $\psi(s) = 3s^2 - 2s^3$ for $0 \leq s \leq 1$, $\psi = 0$ for $s < 0$ and $\psi = 1$ for $s > 1$, a $C^1$ function. For a ball $B(x, r)$ and $j > 1/r$ let
$\varphi_j(y) = \psi(j(r - \|y - x\|))$; it is $C^1$, since it is constant near $y = x$, takes values in $[0, 1]$, vanishes outside $B(x, r)$, and $\varphi_j \to \mathbf{1}_{B(x, r)}$
pointwise. By dominated convergence $\int_{B(x, r)}g = 0$ for all balls, and by the Lebesgue differentiation theorem (Theorem~\ref{measure-theory:theorem-lebesgue-differentiation}),
$g(x) = \lim_{r \to 0}A_rg(x) = 0$ for almost every $x$.
\end{proof}

\begin{theorem}[Rademacher]\label{measure-theory:theorem-rademacher}
Every Lipschitz map $f \colon \RR^n \to \RR^m$ is differentiable at almost every point of $\RR^n$.
\end{theorem}

\begin{proof}
It suffices to treat $m = 1$, applying the result to the components. Let $L$ be a Lipschitz constant. For a unit vector $v$ and $x \in \RR^n$ let
$D_vf(x) = \lim_{t \to 0}(f(x + tv) - f(x))/t$ where the limit exists.

\emph{Step 1: for each $v$, $D_vf(x)$ exists for almost every $x$.} As the difference quotients are continuous in $t \neq 0$, the limit exists if and only if the $\limsup$ and $\liminf$ over
rational $t \to 0$ are equal and finite, so the set $N_v$ where it fails is Borel, and $D_vf$ is Borel measurable where defined, with $|D_vf| \leq L$. Choose an orthogonal map $R$ with $Re_1 = v$;
it preserves $\lambda$ (Lemma~\ref{measure-theory:lemma-linear-change}), and $f \circ R$ is Lipschitz. For each $y \in \RR^{n-1}$, $t \mapsto f(R(t, y))$ is a Lipschitz function of one variable,
differentiable almost everywhere (Example~\ref{measure-theory:example-absolutely-continuous}); so every line parallel to $e_1$ meets $R^{-1}(N_v)$ in a null set, and by Tonelli's theorem
(Theorem~\ref{measure-theory:theorem-fubini}) $\lambda(N_v) = \lambda(R^{-1}(N_v)) = 0$. In particular the gradient $\nabla f = (D_{e_1}f, \ldots, D_{e_n}f)$ exists almost everywhere, and
$|\nabla f| \leq \sqrt nL$.

\emph{Step 2: $D_vf = \langle v, \nabla f\rangle$ almost everywhere.} Let $\varphi$ be $C^1$ with compact support. By translation invariance,
\[
\int\frac{f(x + tv) - f(x)}{t}\varphi(x)\,dx = -\int f(x)\frac{\varphi(x) - \varphi(x - tv)}{t}\,dx .
\]
Let $t = 1/j \to 0$. On the left the integrands converge almost everywhere to $D_vf\,\varphi$ and are bounded by $L|\varphi|$; on the right they converge to $f\,D_v\varphi = f\langle v, \nabla\varphi\rangle$
and are bounded by a constant times the indicator of a compact set, by the mean value inequality and the continuity of $f$. By dominated convergence,
$\int D_vf\,\varphi = -\int f\langle v, \nabla\varphi\rangle = -\sum_iv_i\int f\,\partial_i\varphi = \sum_iv_i\int D_{e_i}f\,\varphi$, the last equality being the first one for $v = e_i$. So
$\int(D_vf - \langle v, \nabla f\rangle)\varphi = 0$ for all $\varphi$, and Lemma~\ref{measure-theory:lemma-vanishing-integrals} applies.

\emph{Step 3: differentiability.} Let $(v_k)$ be a dense sequence in the unit sphere, and let $A$ be the set of points where $\nabla f$ and all $D_{v_k}f$ exist and
$D_{v_k}f = \langle v_k, \nabla f\rangle$; its complement is a countable union of null sets. Fix $x \in A$, and for a unit vector $v$ and $t \neq 0$ put
$Q(v, t) = (f(x + tv) - f(x))/t - \langle v, \nabla f(x)\rangle$. Then $|Q(v, t) - Q(v^{\prime}, t)| \leq (L + \sqrt nL)|v - v^{\prime}|$. Given $\varepsilon > 0$, by compactness of the
sphere there is $N$ such that every unit vector is within $\varepsilon/(2(1 + \sqrt n)L)$ of some $v_k$ with $k \leq N$, and there is $\delta > 0$ with $|Q(v_k, t)| < \varepsilon/2$ for
$0 < |t| < \delta$ and $k \leq N$. Then $|Q(v, t)| < \varepsilon$ for all unit vectors $v$ and $0 < |t| < \delta$, that is, $|f(x + h) - f(x) - \langle\nabla f(x), h\rangle| \leq \varepsilon|h|$ for
$|h| < \delta$. So $f$ is differentiable at $x$, with derivative $h \mapsto \langle\nabla f(x), h\rangle$.
\end{proof}

\begin{counterexample}\label{measure-theory:counterexample-rademacher}
\begin{enumerate}
\item Continuity is not enough: the Takagi function (Theorem~\ref{real-analysis:theorem-takagi}) is continuous and differentiable nowhere, so it is not Lipschitz on any interval.
\item The exceptional null set can be dense. Let $q_1, q_2, \ldots$ be an enumeration of $\QQ$ and $f(x) = \sum_j2^{-j}|x - q_j|$, a $1$-Lipschitz function on $\RR$. For $h > 0$,
\[
\frac{f(q_k + h) - f(q_k)}{h} - \frac{f(q_k) - f(q_k - h)}{h} = \sum_j2^{-j}\,\frac{|q_k - q_j + h| + |q_k - q_j - h| - 2|q_k - q_j|}{h} \geq 2^{1-k},
\]
since every term is $\geq 0$ by the triangle inequality and the term $j = k$ is $2^{1-k}$. If $f$ were differentiable at $q_k$ the left side would tend to $0$; so $f$ is differentiable at no
rational point.
\item Existence of all directional derivatives at a point does not imply differentiability there, even for Lipschitz functions: $f(x, y) = x^2y/(x^2 + y^2)$, $f(0, 0) = 0$, is Lipschitz,
has $D_vf(0) = f(v)$ for every unit vector $v$, and $v \mapsto f(v)$ is not linear. Step 3 of the proof is where the Lipschitz condition turns directional derivatives into a derivative.
\end{enumerate}
\end{counterexample}

\section{The area formula}\label{measure-theory:section-area-formula}

For a linear map $T \colon \RR^n \to \RR^n$, $\lambda^n(T(A)) = |\det T|\lambda^n(A)$. For a linear map $L \colon \RR^k \to \RR^m$ with $k \leq m$ the image of a cube is a
$k$-dimensional parallelepiped in $\RR^m$, and its $k$-dimensional volume is measured by $\mathcal{H}^k$. The factor replacing $|\det T|$ is the Jacobian of $L$.

\begin{definition}\label{measure-theory:definition-jacobian}
The \emph{Jacobian} of a linear map $L \colon \RR^k \to \RR^m$, $k \leq m$, is $JL = \sqrt{\det(L^TL)}$. For a map $f$ differentiable at $x$, $Jf(x) = J(Df(x))$.
\end{definition}

\begin{lemma}\label{measure-theory:lemma-jacobian}
Let $L \colon \RR^k \to \RR^m$ be linear, $k \leq m$.
\begin{enumerate}
\item $(JL)^2$ is the sum of the squares of the $k \times k$ minors of the matrix of $L$, and $JL = 0$ if and only if $L$ is not injective. For $k = m$, $JL = |\det L|$.
\item If $L$ is injective, $L = O \circ S$ with $S = (L^TL)^{1/2}$ positive definite symmetric on $\RR^k$, $\det S = JL$, and $O \colon \RR^k \to \RR^m$ an isometry onto $L(\RR^k)$.
\item $\mathcal{H}^k(L(A)) = JL\,\lambda^k(A)$ for every Lebesgue measurable $A \subseteq \RR^k$.
\end{enumerate}
\end{lemma}

\begin{proof}
(1) By Cauchy--Binet (Proposition~\ref{multilinear-algebra:proposition-minors}(2)) applied to $L^T$ and $L$, $\det(L^TL) = \sum_S\det(L_{S})^2$, where $L_S$ is the $k \times k$ submatrix with rows
$S$. By Proposition~\ref{multilinear-algebra:proposition-minors}(3), some minor is nonzero if and only if $L$ has rank $k$. For $k = m$, $\det(L^TL) = (\det L)^2$.

(2) $L^TL$ is symmetric with $\langle L^TLv, v\rangle = |Lv|^2 > 0$ for $v \neq 0$, so by the spectral theorem (Theorem~\ref{linear-algebra:theorem-spectral}) it has a positive definite square root $S$,
with $\det S = \sqrt{\det L^TL}$. Then $O = LS^{-1}$ satisfies $O^TO = S^{-1}L^TLS^{-1} = I$, so $|Ov| = |v|$, and $O(\RR^k) = L(\RR^k)$.

(3) If $L$ is not injective, $L(\RR^k)$ lies in a linear subspace $W$ of dimension $k - 1$; an isometry $\RR^{k-1} \to W$ and Theorem~\ref{measure-theory:theorem-hausdorff-lebesgue} give
$\mathcal{H}^{k-1}(W \cap B) < \infty$ for every ball $B$, so $\mathcal{H}^k(W) = 0$ by Proposition~\ref{measure-theory:proposition-hausdorff}(4), and both sides vanish. If $L$ is injective, $O$ is an
isometry from $\RR^k$ onto $L(\RR^k)$, which preserves $\mathcal{H}^k$, so $\mathcal{H}^k(L(A)) = \mathcal{H}^k(S(A)) = \lambda^k(S(A)) = |\det S|\lambda^k(A)$ by
Theorem~\ref{measure-theory:theorem-hausdorff-lebesgue} and Lemma~\ref{measure-theory:lemma-linear-change}; for Lebesgue measurable $A$ write $A = B \cup N$ with $B$ Borel and $N$ contained in a Borel null
set, which $S$ maps to a null set.
\end{proof}

So $JL$ is the $k$-dimensional volume of the parallelepiped spanned by $Le_1, \ldots, Le_k$, the image of the unit cube. For a curve, $J\gamma(t) = |\gamma^{\prime}(t)|$; for a parametrised
surface $x(u, v)$ in $\RR^3$, $Jx = |x_u \times x_v|$, since the three $2 \times 2$ minors of $Dx$ are the components of $x_u \times x_v$ up to sign.

\begin{theorem}[Area formula]\label{measure-theory:theorem-area-formula}
Let $U \subseteq \RR^k$ be open and $f \colon U \to \RR^m$ an injective $C^1$ map whose differential $Df(x)$ is injective at every point. Then for every Lebesgue measurable $A \subseteq U$ the set
$f(A)$ is $\mathcal{H}^k$-measurable, and
\[
\mathcal{H}^k(f(A)) = \int_AJf\,d\lambda^k .
\]
Consequently $\int_{f(U)}g\,d\mathcal{H}^k = \int_Ug(f(x))\,Jf(x)\,d\lambda^k(x)$ for every Borel function $g \geq 0$ on $\RR^m$.
\end{theorem}

\begin{proof}
For an injective linear $L \colon \RR^k \to \RR^m$ let $m(L) = \min_{|v| = 1}|Lv| > 0$. Fix $t > 1$.

\emph{Step 1: a decomposition of $U$.} Call a pair $(B, L)$ of an open ball $B \subseteq U$ with rational centre and radius and an injective linear map $L$ with rational matrix \emph{good} if, for all
$x, y \in B$,
\begin{equation}\label{measure-theory:equation-good-pair}
t^{-1}|L(x - y)| \leq |f(x) - f(y)| \leq t|L(x - y)| \qquad\text{and}\qquad t^{-k}JL \leq Jf(x) \leq t^kJL .
\end{equation}
Every $x_0 \in U$ lies in the ball of some good pair. Indeed, let $\mu_0 = m(Df(x_0)) > 0$ and $\eta = (1 - t^{-1})\mu_0/2$. Since $|m(L) - m(L^{\prime})| \leq \|L - L^{\prime}\|$, $J$ is
continuous, and $Jf(x_0) > 0$, there is $L$ with rational matrix such that $\|Df(x_0) - L\| < \min(\eta, \mu_0)/2$ and $t^{-k}JL < Jf(x_0) < t^kJL$; then $m(L) \geq \mu_0/2$, so
$\eta \leq (1 - t^{-1})m(L)$. By continuity of $Df$ and of $J$ there is a rational ball $B \ni x_0$ on which $\|Df - L\| \leq \eta$ and still $t^{-k}JL \leq Jf \leq t^kJL$. For $x, y \in B$, by the mean value inequality (Proposition~\ref{real-analysis:proposition-derivative-basic}(4))
applied to $f - L$ on the convex set $B$, $|f(x) - f(y) - L(x - y)| \leq \eta|x - y| \leq (1 - t^{-1})|L(x - y)|$, which gives
$t^{-1}|L(x - y)| \leq |f(x) - f(y)| \leq (2 - t^{-1})|L(x - y)| \leq t|L(x - y)|$. Enumerate the countably many good pairs as $(B_j, L_j)$, and let $E_j = B_j \setminus \bigcup_{i < j}B_i$; these are
disjoint Borel sets covering $U$.

\emph{Step 2: estimates on one piece.} Let $A \subseteq E_j$ be Lebesgue measurable. By \eqref{measure-theory:equation-good-pair} the map $L_j(x) \mapsto f(x)$, from $L_j(B_j)$ to $f(B_j)$, is
$t$-Lipschitz, and so is its inverse. By Proposition~\ref{measure-theory:proposition-hausdorff}(3) and Lemma~\ref{measure-theory:lemma-jacobian}(3),
\[
t^{-k}JL_j\,\lambda^k(A) \leq \mathcal{H}^k(f(A)) \leq t^kJL_j\,\lambda^k(A), \qquad t^{-k}JL_j\,\lambda^k(A) \leq \int_AJf \leq t^kJL_j\,\lambda^k(A),
\]
the second pair by \eqref{measure-theory:equation-good-pair}. Hence $t^{-2k}\int_AJf \leq \mathcal{H}^k(f(A)) \leq t^{2k}\int_AJf$; here $\mathcal{H}^k$ is the outer measure.

\emph{Step 3: measurability.} A Lebesgue measurable $A \subseteq U$ is $K \cup N$ with $K$ a countable union of compact sets and $N$ null (Theorem~\ref{measure-theory:theorem-lebesgue-measure}).
Then $f(K)$ is a countable union of compact sets, hence Borel, and $\mathcal{H}^k(f(N)) \leq \sum_j\mathcal{H}^k(f(N \cap E_j)) = 0$ by Step 2. So $f(A)$ differs from a Borel set by an
$\mathcal{H}^k$-null set, and it is $\mathcal{H}^k$-measurable (Theorem~\ref{measure-theory:theorem-caratheodory}).

\emph{Step 4.} For measurable $A \subseteq U$, the sets $f(A \cap E_j)$ are measurable and disjoint, since $f$ is injective. Summing Step 2 over $j$,
$t^{-2k}\int_AJf \leq \mathcal{H}^k(f(A)) \leq t^{2k}\int_AJf$, and $t > 1$ was arbitrary. For a Borel set $G \subseteq \RR^m$, applying this to $A = f^{-1}(G)$ gives the integral formula for
$g = \mathbf{1}_G$; it follows for simple functions by linearity, and for Borel $g \geq 0$ by monotone convergence (Theorem~\ref{measure-theory:theorem-monotone-convergence}).
\end{proof}

For $k = m$ the area formula is the change of variables formula (Theorem~\ref{measure-theory:theorem-change-of-variables}), with $\mathcal{H}^m = \lambda^m$. The proof is the same idea: on small pieces
$f$ is almost linear, and linear maps multiply measure by the Jacobian.

\begin{example}\label{measure-theory:example-area-formula}
\begin{enumerate}
\item \emph{Length of curves.} For an injective $C^1$ curve $\gamma \colon (a, b) \to \RR^m$ with $\gamma^{\prime} \neq 0$, $\mathcal{H}^1(\gamma((a, b))) = \int_a^b|\gamma^{\prime}(t)|\,dt$. So the length
of a curve, defined by parametrisations, is the one-dimensional Hausdorff measure of its image, which depends only on the set.
\item \emph{Area of graphs.} For $h \colon U \to \RR$ of class $C^1$, $U \subseteq \RR^k$ open, the graph map $f(x) = (x, h(x))$ is injective with injective differential, and
$Jf = \sqrt{\det(I + \nabla h\,\nabla h^T)} = \sqrt{1 + |\nabla h|^2}$, since $I + ww^T$ is the identity on $w^\perp$ and multiplies $w$ by $1 + |w|^2$. So the graph has area
$\int_U\sqrt{1 + |\nabla h|^2}\,d\lambda^k$.
\item \emph{The area of the sphere.} The upper hemisphere of $S^2$ is the graph of $h(x) = \sqrt{1 - |x|^2}$ over the open unit disc, with $|\nabla h|^2 = |x|^2/(1 - |x|^2)$ and
$Jf = (1 - |x|^2)^{-1/2}$. In polar coordinates (Example~\ref{measure-theory:example-gaussian}), $\mathcal{H}^2 = \int_0^{2\pi}\int_0^1 r(1 - r^2)^{-1/2}\,dr\,d\theta = 2\pi$. The equator has finite $\mathcal{H}^1$-measure by (1), so it is
$\mathcal{H}^2$-null (Proposition~\ref{measure-theory:proposition-hausdorff}(4)), and $\mathcal{H}^2(S^2) = 4\pi$.
\end{enumerate}
\end{example}

\begin{counterexample}\label{measure-theory:counterexample-area-injective}
Injectivity is needed: for $f(t) = (\cos t, \sin t)$ on $(0, 4\pi)$, $\int Jf = 4\pi$, while the image, the circle minus one point, has $\mathcal{H}^1$-measure $2\pi$. The integral counts the points of the
image with multiplicity: in general $\int_AJf = \int_{\RR^m}\#(A \cap f^{-1}(y))\,d\mathcal{H}^k(y)$ (Remark~\ref{measure-theory:remark-area-coarea}).
\end{counterexample}

\begin{remark}\label{measure-theory:remark-area-coarea}
\emph{The general area and coarea formulas.} By Rademacher's theorem a Lipschitz map $f \colon \RR^k \to \RR^m$ has a Jacobian almost everywhere, and the area formula holds in the generality
\[
\int_AJf\,d\lambda^k = \int_{\RR^m}\#\bigl(A \cap f^{-1}(y)\bigr)\,d\mathcal{H}^k(y) \qquad (k \leq m),
\]
for Lebesgue measurable $A$, where $\#$ is the number of points \cite[Section 3.3]{evans-gariepy}, \cite[Section 3.2]{federer-gmt}. For $k \geq m$ there is a dual statement, the \emph{coarea
formula}: with $Jf = \sqrt{\det(Df\,Df^T)}$,
\[
\int_AJf\,d\lambda^k = \int_{\RR^m}\mathcal{H}^{k-m}\bigl(A \cap f^{-1}(y)\bigr)\,d\lambda^m(y) \qquad (k \geq m),
\]
which slices $A$ by the level sets of $f$ instead of parametrising it \cite[Section 3.4]{evans-gariepy}. For $f(x) = |x|$ on $\RR^n$, $Jf = |\nabla f| = 1$ away from $0$, and the coarea formula
gives integration in spherical shells, $\int_{\RR^n}g\,d\lambda = \int_0^\infty\bigl(\int_{|x| = r}g\,d\mathcal{H}^{n-1}\bigr)dr$; for a $C^1$ function $u$ it gives $\int|\nabla u| =
\int_\RR\mathcal{H}^{n-1}(\{u = s\})\,ds$, the starting point of the isoperimetric and Sobolev inequalities. A set is \emph{$k$-rectifiable} if it is covered, up to an $\mathcal{H}^k$-null set,
by countably many Lipschitz images of subsets of $\RR^k$; these sets, which include all $k$-dimensional submanifolds, are the objects on which geometric measure theory does calculus, and
the area formula is the bridge between them and Lebesgue measure.
\end{remark}

\section{Exercises}\label{measure-theory:section-exercises}

\begin{exercise}\label{measure-theory:exercise-cantor}
Show that the Cantor set, obtained from $[0, 1]$ by repeatedly removing open middle thirds, is compact, uncountable and of Lebesgue measure zero.
\end{exercise}

\begin{solution}
It is an intersection of closed subsets of $[0, 1]$, hence compact. After $k$ steps it is covered by $2^k$ intervals of length $3^{-k}$, so its measure is at most $(2/3)^k$. Its points are the numbers
$\sum_k a_k3^{-k}$ with $a_k \in \{0, 2\}$, and distinct sequences give distinct numbers, so it is in bijection with $\{0, 1\}^{\NN}$, which is uncountable (Theorem~\ref{sets:theorem-cantor}).
\end{solution}

\begin{exercise}\label{measure-theory:exercise-borel-cantelli}
(Borel--Cantelli.) If $\sum_k\mu(A_k) < \infty$, show that almost every point lies in only finitely many $A_k$.
\end{exercise}

\begin{solution}
By monotone convergence $\int\sum_k\mathbf{1}_{A_k} = \sum_k\mu(A_k) < \infty$, so $\sum_k\mathbf{1}_{A_k} < \infty$ a.e.
\end{solution}

\begin{exercise}\label{measure-theory:exercise-no-half-density}
Show that there is no Lebesgue measurable $A \subseteq \RR$ with $\lambda(A \cap I) = \frac12\lambda(I)$ for every bounded interval $I$.
\end{exercise}

\begin{solution}
For such $A$, $\lambda(A \cap B(x, r))/\lambda(B(x, r)) = \frac12$ for all $x$ and $r$. By Corollary~\ref{measure-theory:corollary-density} this ratio tends to $1$ at
almost every point of $A$ and to $0$ at almost every point outside $A$; so $A$ and $\RR \setminus A$ are both null, which is impossible.
\end{solution}

\begin{exercise}\label{measure-theory:exercise-maximal-not-integrable}
Show that $Mf \notin L^1(\RR^n)$ for every $f \in L^1(\RR^n)$ that is not zero almost everywhere; so the maximal inequality cannot be improved to
$\|Mf\|_1 \leq C\|f\|_1$.
\end{exercise}

\begin{solution}
Choose $R$ with $a = \int_{B(0, R)}|f| > 0$. For $\|x\| \geq R$, the ball $B(x, 2\|x\|)$ contains $B(0, R)$, so
$Mf(x) \geq a/\lambda(B(x, 2\|x\|)) = a/(2^n\|x\|^n\lambda(B(0, 1)))$. The function $\|x\|^{-n}$ is not integrable on $\{\|x\| \geq R\}$: the shells
$2^kR \leq \|x\| < 2^{k+1}R$ have measure $(2^n - 1)2^{kn}R^n\lambda(B(0, 1))$, and on each of them $\|x\|^{-n} \geq 2^{-(k+1)n}R^{-n}$, so each contributes at least the
positive constant $(1 - 2^{-n})\lambda(B(0, 1))$ to the integral.
\end{solution}

\begin{exercise}\label{measure-theory:exercise-lp-inclusions}
Show that $L^q(\mu) \subseteq L^p(\mu)$ for $p \leq q$ if $\mu(X) < \infty$, and that there is no inclusion in either direction for Lebesgue measure on $\RR$.
\end{exercise}

\begin{solution}
By H\"older with exponents $q/p$ and its conjugate, $\int|f|^p \leq (\int|f|^q)^{p/q}\mu(X)^{1 - p/q}$. On $\RR$, $|x|^{-1/2}\mathbf{1}_{[-1,1]}$ is in $L^1$ but not $L^2$, and $(1 + |x|)^{-1}$ is in $L^2$ but not $L^1$.
\end{solution}

\begin{exercise}\label{measure-theory:exercise-convolution}
For $f \in L^1(\RR^n)$ and $g \in L^p(\RR^n)$, show using Fubini's theorem that $(f * g)(x) = \int f(x - y)g(y)\,dy$ is defined a.e. and $\|f * g\|_p \leq \|f\|_1\|g\|_p$.
\end{exercise}

\begin{exercise}\label{measure-theory:exercise-ball-volume}
Show that the volume of the unit ball in $\RR^n$ is $\omega_n = \pi^{n/2}/\Gamma(\frac n2 + 1)$, where $\Gamma(s) = \int_0^\infty t^{s-1}e^{-t}\,dt$, by computing $\int_{\RR^n}e^{-\|x\|^2}$ in two ways.
\end{exercise}

\begin{exercise}\label{measure-theory:exercise-middle-half-cantor}
Let $K \subseteq [0, 1]$ be the set obtained by removing the open middle half $(\frac14, \frac34)$ of $[0, 1]$, then the open middle halves of the two remaining intervals, and so on; equivalently, the
numbers $\sum_{k \geq 1}a_k4^{-k}$ with all $a_k \in \{0, 3\}$. Show that $\dim_{\mathcal{H}}K = \frac12$.
\end{exercise}

\begin{solution}
Follow Example~\ref{measure-theory:example-cantor-dimension} with $3$ replaced by $4$ and the digits $\{0, 2\}$ by $\{0, 3\}$; now $s = \log 2/\log 4 = \frac12$, so $4^s = 2$. The $2^k$ intervals of
generation $k$ have length $4^{-k}$, which gives $\mathcal{H}^{1/2}(K) \leq \alpha(\frac12)2^k4^{-k/2} = \alpha(\frac12)$. Two of them first differing at position $j \leq k$ have left ends differing by at least
$3 \cdot 4^{-j} - \sum_{j < i \leq k}3 \cdot 4^{-i} \geq 2 \cdot 4^{-j}$, so they are at distance at least $4^{-k}$. The function $F(\sum_ka_k4^{-k}) = \sum_k(a_k/3)2^{-k}$ on $K$, extended to
$[0, 1]$ as in Counterexample~\ref{measure-theory:counterexample-cantor-function}, is continuous and nondecreasing, and its measure $\mu_F$ (Lemma~\ref{measure-theory:lemma-stieltjes}) gives mass
$2^{-k}$ to each interval of generation $k$ and mass $0$ to $[0, 1] \setminus K$. A closed set of diameter $d$ with $4^{-k-1} \leq d < 4^{-k}$ meets at most one interval of generation $k$, so its
$\mu_F$-measure is at most $2^{-k} = 2 \cdot 4^{-(k+1)/2} \leq 2d^{1/2}$, and as in the example $\mathcal{H}^{1/2}(K) \geq \alpha(\frac12)/2 > 0$.
\end{solution}

\begin{exercise}\label{measure-theory:exercise-lipschitz-graph}
Let $h \colon \RR^n \to \RR$ be Lipschitz with constant $L$ and $U \subseteq \RR^n$ a bounded open set. Show that the graph $\Gamma = \{(x, h(x)) : x \in U\}$ satisfies
$\lambda^n(U) \leq \mathcal{H}^n(\Gamma) \leq (1 + L^2)^{n/2}\lambda^n(U)$, so that $\dim_{\mathcal{H}}\Gamma = n$, even though $h$ need not be differentiable everywhere.
\end{exercise}

\begin{solution}
The map $g(x) = (x, h(x))$ satisfies $|g(x) - g(y)|^2 = |x - y|^2 + |h(x) - h(y)|^2 \leq (1 + L^2)|x - y|^2$, so it is $\sqrt{1 + L^2}$-Lipschitz, and its inverse, the projection
$\Gamma \to U$, is $1$-Lipschitz. By Proposition~\ref{measure-theory:proposition-hausdorff}(3) and Theorem~\ref{measure-theory:theorem-hausdorff-lebesgue},
$\lambda^n(U) = \mathcal{H}^n(U) \leq \mathcal{H}^n(\Gamma) \leq (1 + L^2)^{n/2}\mathcal{H}^n(U)$. Both bounds are positive and finite, so $\mathcal{H}^s(\Gamma)$ is $\infty$ for $s < n$ and $0$ for
$s > n$ by Proposition~\ref{measure-theory:proposition-hausdorff}(4). (By the general area formula of Remark~\ref{measure-theory:remark-area-coarea} and Rademacher's theorem,
$\mathcal{H}^n(\Gamma) = \int_U\sqrt{1 + |\nabla h|^2}$.)
\end{solution}

\begin{exercise}\label{measure-theory:exercise-torus-area}
For $R > r > 0$ let $x(u, v) = ((R + r\cos u)\cos v, (R + r\cos u)\sin v, r\sin u)$. Show that $x$ is injective with injective differential on $(0, 2\pi)^2$, that $Jx = r(R + r\cos u)$, and that
the torus $T = x(\RR^2)$ has $\mathcal{H}^2(T) = 4\pi^2Rr$.
\end{exercise}

\begin{solution}
$x_u = (-r\sin u\cos v, -r\sin u\sin v, r\cos u)$ and $x_v = (-(R + r\cos u)\sin v, (R + r\cos u)\cos v, 0)$ are orthogonal, of lengths $r$ and $R + r\cos u > 0$, so
$Jx = |x_u \times x_v| = r(R + r\cos u)$ and $Dx$ is injective. The point $x(u, v)$ determines $v$ modulo $2\pi$ (its projection to the plane $z = 0$ has argument $v$) and then $u$ modulo
$2\pi$, so $x$ is injective on $(0, 2\pi)^2$. Its image is $T$ minus the two circles $u = 0$ and $v = 0$, which have finite $\mathcal{H}^1$-measure and are therefore $\mathcal{H}^2$-null. By the area
formula, $\mathcal{H}^2(T) = \int_0^{2\pi}\int_0^{2\pi}r(R + r\cos u)\,du\,dv = 4\pi^2Rr$.
\end{solution}
